% The S6 record and the integral monodromy of its (3,4,infinity) period system (version 3).
% Prepared by Claude (Anthropic), model claude-opus-5-5 (Opus 5.5).
% Compile with lualatex (three times).
% The S6 record and the integral monodromy of its (3,4,infinity) period system (version 3).
% Prepared by Claude (Anthropic), model claude-opus-5-5 (Opus 5.5).
% Compile with lualatex (twice).
\documentclass[11pt,a4paper]{article}
\usepackage[a4paper,margin=24mm,headheight=14pt]{geometry}
\usepackage{fontspec}
\directlua{luaotfload.add_fallback("readerfallback",{"TeX Gyre Pagella Math:mode=harf;","STIX:mode=harf;","DejaVu Serif:mode=harf;","FreeSerif:mode=harf;"})}
\setmainfont{TeX Gyre Pagella}[RawFeature={fallback=readerfallback}]
\setsansfont{TeX Gyre Heros}[Scale=MatchLowercase]
\setmonofont{DejaVu Sans Mono}[Scale=0.84]
\usepackage{amsmath,amsthm,mathtools}
\usepackage{graphicx}
\usepackage{unicode-math}
\setmathfont{TeX Gyre Pagella Math}
\usepackage{microtype}
\usepackage{booktabs,longtable,array,tabularx}
\usepackage{enumitem}
\setlist{itemsep=2pt,topsep=4pt,leftmargin=*}
\usepackage[dvipsnames]{xcolor}
\usepackage{fancyhdr}
\usepackage{titlesec}
\usepackage[hidelinks,bookmarksnumbered]{hyperref}
\hypersetup{pdftitle={The S6 record and the integral monodromy of its (3,4,infinity) period system},
  pdfauthor={Claude (Anthropic), model claude-opus-5-5; record by KokunoYumeto},
  pdfsubject={The S6 verification and reconstruction record (Zenodo 10.5281/zenodo.22678442): results with proofs, and the integral monodromy of the period system},
  pdfkeywords={complex structure on the six-sphere, compact complex threefolds, algebraic dimension, torus fibrations, monodromy, paramodular group, Heisenberg group, arithmetic groups}}
\usepackage{bookmark}

\titleformat{\section}{\normalfont\Large\bfseries}{\thesection}{0.8em}{}
\titleformat{\subsection}{\normalfont\large\bfseries}{\thesubsection}{0.7em}{}
\titlespacing*{\section}{0pt}{2.2ex plus 1ex}{1.2ex}
\titlespacing*{\subsection}{0pt}{1.8ex plus .8ex}{0.9ex}

\pagestyle{fancy}
\fancyhf{}
\fancyhead[L]{\small\itshape The $S^6$ record and the integral monodromy of its period system}
\fancyhead[R]{\small\thepage}
\renewcommand{\headrulewidth}{0.3pt}
\fancypagestyle{plain}{\fancyhf{}\fancyfoot[C]{\small\thepage}\renewcommand{\headrulewidth}{0pt}}

\theoremstyle{plain}
\newtheorem{theorem}{Theorem}[section]
\newtheorem{proposition}[theorem]{Proposition}
\newtheorem{lemma}[theorem]{Lemma}
\newtheorem{corollary}[theorem]{Corollary}
\newtheorem{introthm}{Theorem}
\renewcommand{\theintrothm}{\Roman{introthm}}
\theoremstyle{definition}
\newtheorem{definition}[theorem]{Definition}
\newtheorem{example}[theorem]{Example}
\theoremstyle{remark}
\newtheorem{remark}[theorem]{Remark}

% status line printed under a statement
\newcommand{\status}[1]{\par\smallskip\noindent{\small\color{black!70}\textsf{Status.}\ #1}\par\smallskip}
\newcommand{\negres}[2]{\par\smallskip\noindent\textbf{#1}\ #2\par}

% notation
\newcommand{\B}{\mathcal{B}}
\newcommand{\Zs}{\mathcal{Z}}
\newcommand{\C}{\mathbb{C}}
\newcommand{\R}{\mathbb{R}}
\newcommand{\Q}{\mathbb{Q}}
\newcommand{\ZZ}{\mathbb{Z}}
\newcommand{\NN}{\mathbb{N}}
\newcommand{\Ff}{\mathbb{F}}
\newcommand{\Pp}{\mathbb{P}}
\newcommand{\re}{\operatorname{Re}}
\newcommand{\im}{\operatorname{Im}}
\newcommand{\Res}{\operatorname*{Res}}
\newcommand{\sgn}{\operatorname{sgn}}
\newcommand{\Gr}{\operatorname{Gr}}
\newcommand{\cl}{\overline}
\newcommand{\brkus}{\textunderscore\allowbreak}
\newcommand{\note}[1]{{\ttfamily\let\_\brkus #1}}
\newcommand{\fn}[1]{{\ttfamily\let\_\brkus #1}}
\newcommand{\lbl}[1]{\textsf{#1}}

\setlength{\parskip}{3pt}
\setlength{\emergencystretch}{2.5em}
\setlength{\parindent}{0pt}

\newcommand{\T}{\mathcal{T}}
\newcommand{\Hs}{\mathcal{H}}
\newcommand{\lcm}{\operatorname{lcm}}
\newcommand{\TV}{\operatorname{TV}}
\newcommand{\tr}{\operatorname{tr}}
\newcommand{\Var}{\operatorname{Var}}


\begin{document}
\thispagestyle{plain}
\begin{center}
{\LARGE\bfseries The \texorpdfstring{$S^6$}{S6} record and the integral monodromy\\[2pt] of its \texorpdfstring{$(3,4,\infty)$}{(3,4,infinity)} period system}\\[8pt]
{\Large Its results, with proofs, a correction to a published step,\\ and its bridges}\\[14pt]
{\large A paper prepared from an audit of the record\\ by Claude (Anthropic), model \texttt{claude-opus-5-5}}\\[8pt]
{\normalsize Record: KokunoYumeto (Zenodo 10.5281/zenodo.22678442); the construction is that of the manuscript circulated by Levent Alpöge}\\[4pt]
{\normalsize Audit 26--27 September 2026; version~3, 28 September 2026 (the text of version~2 with its framing corrected, and the integral monodromy of the period system; Appendix~\ref{app:changes})}\\[4pt]
{\small Published at the author's request ahead of the author's review; the author may revise or withdraw it.}\\[2pt]
{\small Companion to the paper on the split-zero programme around the Riemann zeta function, version~3, in the Zenodo record with concept DOI \href{https://doi.org/10.5281/zenodo.22678085}{10.5281/zenodo.22678085}}
\end{center}

\begin{abstract}
\noindent
Whether the six-sphere carries a complex structure is a classical open question. A manuscript circulated by Levent Alpöge claims one: a family of complex two-dimensional tori over the projective line, with local monodromies $T_1,T_2$ of orders $3$ and $4$, completed at its three special fibres, whose total space would be a compact complex threefold diffeomorphic to $S^6$ of algebraic dimension one. Such a threefold would contradict a theorem of Campana, Demailly and Peternell. The Zenodo record 10.5281/zenodo.22678442 verifies and reconstructs the manuscript. This paper sets out what the record establishes, with proofs where they are given here, and adds new results.

Two parts of the construction are determined completely here. The first is the step of the published proof that the record's correction note isolates. The paper audits the note and proves that the torsion responsible for the failure of the step is present exactly at the non-normal points of the fibre, and that at a reduced fibre with normal crossings the group which the step asserts to vanish is the space of sections of a line bundle on the normalization. The second is the integral period system. Its monodromy group is the kernel of an explicit character $\Phi$ of order $12$ on the stabiliser of an isotropic vector: a normal subgroup of index $12$ with cyclic quotient, a congruence subgroup of exact level $12$, arithmetic in its Zariski closure, and a non-split extension of $SL_2(\ZZ)$ by a Heisenberg group. The Levi half of $\Phi$ is, up to sign, the peripheral character of the record's higher-torus update. Its Heisenberg half is defined by the odd quadratic refinement of the mod-$2$ form, the odd spin structure of a torus, and on the standard Heisenberg subgroup it is the Heisenberg multiplier of the odd Jacobi theta function. The orbits of the monodromy group on the affine hyperplanes used by the Erdős--Straus workbench are classified at every level, and the exact value set of an octonionic cubic of the record's key advances is determined.

The paper also describes the bridges from the period system to paramodular forms, spin structures, degenerations of abelian surfaces and the author's other workbenches, with what is established about each. The existence of the threefold and the counterexample claim are not verified here.
\end{abstract}

\section*{Introduction}\phantomsection\label{sec:about}
\addcontentsline{toc}{section}{Introduction}

\subsection*{The problem and this record}

Campana, Demailly and Peternell proved that a compact complex threefold homeomorphic to $S^6$ would have algebraic dimension $0$ \cite[Corollary~2.3]{CDP20}. A complex structure on $S^6$ that fibres holomorphically over $\Pp^1$ would have algebraic dimension at least one, and would contradict that theorem. The manuscript \cite{AlpogeS6} claims such a fibration, built from the $(3,4,\infty)$ triangle group and a representation of it on $\ZZ^4$ (Section~\ref{sec:candidate}). The record, created by KokunoYumeto, with a verification by Codex and two review records by ChatGPT Pro, reconstructs the manuscript, repairs its local defects, and labels the construction and the counterexample as verified (Sections~\ref{sec:record} and~\ref{sec:ledger}). Its correction note isolates a step in the published proof that is not a formal consequence of the sequences displayed there; the manuscript had located the failure of that step at its own cusp fibre (Section~\ref{sec:correction}). Two later supplements build on the construction's data (Section~\ref{sec:later}). Philip Engel's exposition \cite{EngelS6} presents the construction with the stated aim of a self-contained proof.

This paper states what the record claims, proves what can be proved from the record's local, finite and algebraic statements, and adds the results below. It does not verify the global construction.

\subsection*{Main results}

The first result concerns the step. Let $S=\{g=0\}$ be a reduced hypersurface in a complex manifold $X$, with conormal sheaf $N^*$ and normalization $\eta:\tilde S\to S$; let $T_S$ be the torsion subsheaf of $\Omega^1_S$ and $\widetilde\Omega^1_S=\Omega^1_S/T_S$. The step asserts that $H^0(S,N^*\otimes A)=0$ and $H^0(S,\widetilde\Omega^1_S\otimes A)=0$ imply $H^0(S,\Omega^1_X|_S\otimes A)=0$.

\begin{introthm}[Propositions~\ref{prop:normal}, \ref{prop:kernel} and~\ref{prop:sharp}]
\textup{(i)} $T_S$ vanishes at $x$ if and only if $S$ is normal at $x$; so the step holds on normal fibres, and on a non-normal fibre it holds exactly when a connecting map $\delta_A$ is injective.
\textup{(ii)} If $S$ has normal crossings and $H^0(S,\widetilde\Omega^1_S\otimes A)=0$, then $H^0(S,\Omega^1_X|_S\otimes A)\cong H^0(\tilde S,\eta^*(N^*\otimes A))$. On a fibre, the conclusion of the step holds if and only if $H^0(\tilde S,\eta^*A)=0$.
\textup{(iii)} In the correction note's test family this group is one-dimensional for every value of the gluing parameter.
\end{introthm}

The normal case of (i) is the manuscript's Remark~10.4; the converse, (ii) and (iii) are new. So the step needs an extra hypothesis exactly at non-normal fibres, and (ii) states that hypothesis for the fibres with normal crossings that the construction uses.

The remaining results concern the integral period system. On $V=\ZZ^4$ with basis $(\gamma,u,w,\delta)$ the local monodromies $T_1,T_2$ fix $\gamma$ and preserve an alternating form $Q_0$ of type $(1,6)$. The stabiliser $G_{\ZZ}$ of $\gamma$ in $Sp(Q_0,\ZZ)$ is a semidirect product $SL_2(\ZZ)\ltimes\mathrm{Heis}(\ZZ)$, whose elements we write $L_Mh(p,r,s)$ (Lemma~\ref{lem:stab}). Let $\operatorname{ab}:SL_2(\ZZ)\to\ZZ/12$ be the abelianisation with $\operatorname{ab}\begin{psmallmatrix}1&1\\0&1\end{psmallmatrix}=1$, and put
\[
\Phi\bigl(L_Mh(p,r,s)\bigr)=\operatorname{ab}(M)+s-6(p^2+pr+r^2)\bmod12 .
\]

\begin{introthm}[Theorem~\ref{thm:mono}, Proposition~\ref{prop:twist}]
$\Phi$ is a surjective homomorphism, and the monodromy group $G=\langle T_1,T_2\rangle$ is its kernel. So $G$ is normal of index $12$ in $G_{\ZZ}$, with cyclic quotient generated by the central transvection $\delta\mapsto\delta+\gamma$. It maps onto $SL_2(\ZZ)$, it contains the principal congruence subgroup of level $12$ and none of level $4$ or $6$, and it is Zariski dense and arithmetic in the algebraic stabiliser of $\gamma$. The extension $1\to G\cap\mathrm{Heis}(\ZZ)\to G\to SL_2(\ZZ)\to1$ does not split.
\end{introthm}

\begin{introthm}[Proposition~\ref{prop:heis}, Theorem~\ref{thm:periph}]
\textup{(i)} The Heisenberg part $G\cap\mathrm{Heis}(\ZZ)$ is cut out by the odd quadratic refinement $p^2+pr+r^2$ of the mod-$2$ symplectic form, the only refinement invariant under $SL_2(\ZZ)$. On the standard Heisenberg subgroup it is the kernel of the character $(-1)^{p+r+pr+\kappa}$, the Heisenberg multiplier of the odd Jacobi theta function.
\textup{(ii)} The marked peripheral character $\chi$ of the record's higher-torus update equals $-\operatorname{ab}$ on the Levi parts, its values $(-1,4,-3)$ are those of $\operatorname{ab}$ on the Levi parts of $T_0,T_1,T_2$, and on $G$ it is determined by the Heisenberg coordinates: $\chi(M)\equiv s-6(p^2+pr+r^2)\pmod{12}$.
\end{introthm}

\begin{introthm}[Theorem~\ref{thm:orbitsS6}]
For every $D\ge1$, the orbits of $G$ on the affine hyperplane $H_D$ of the dual lattice modulo $D$ are the sets on which $\gcd(b,c,\gcd(D,6))$ is constant, where $b,c$ are the $\hat u$- and $\hat w$-coordinates, and they coincide with the orbits of $G_{\ZZ}$.
\end{introthm}

These results are new with respect to the record's files that were searched (\fn{00\_}--\fn{03\_}, \fn{14\_}--\fn{16\_} and \fn{26\_}): they use the period system, but contain no statement on the index, level, Zariski closure or arithmeticity of its monodromy group. The archive \fn{28\_} was not searched. The paper also proves the exact value set of the octonionic cubic of the key advance S6-5 (Proposition~\ref{prop:s65}) and a general form of the note's local computation of relative differentials (Proposition~\ref{prop:torsion}). By the latter, the record's correction to the manuscript's Remark~1.3 is right for relative one-forms, while relative two-forms do have torsion along the double curves of the cusp fibre.

\subsection*{What the results mean}

My assessment of the results on the step is that they settle its local content. The step is valid on normal fibres and needs an extra hypothesis on non-normal ones, and at the fibres with normal crossings that hypothesis is the vanishing of $H^0(\tilde S,\eta^*A)$. The global question, whether the hypotheses of \cite{CDP20} force this vanishing, is left open by the note. The record's answer at its own candidate is negative, and that answer depends on the construction, which is not verified here (Question~\ref{q:global}).

My assessment of the monodromy results is that they identify the peripheral character of the record's higher-torus update as the invariant that separates the monodromy group from its arithmetic hull, for the following reasons. The group is as large as it can be subject to one condition, $\Phi=0$. The Levi half of $\Phi$ is, up to sign, the update's character, and the Heisenberg half is the odd theta multiplier. So $G$ is exactly the locus where these two characters agree (Theorem~III(ii)). Because $G$ is arithmetic and not thin, the methods for congruence subgroups apply, and they give the orbit classification that the Erdős--Straus covers need at every level.

\subsection*{The bridges}

In the author's view the connections of this research to other areas are among its most valuable outcomes. The period system connects to several established theories, and Section~\ref{sec:overlaps} treats each with what is established and what is open:
\begin{itemize}
\item \emph{paramodular forms}: $Sp(Q_0,\ZZ)$ is the paramodular group of level $6$ of Gritsenko and Hulek, and $\Phi$ has the shape of one of their characters on it;
\item \emph{spin structures and theta characteristics}: the refinement that cuts out the Heisenberg part is the odd spin structure, which is also the Arf-one spin form of the update's Lie-framed torus;
\item \emph{degenerations}: the cusp monodromy has the Jordan type of Kulikov's type III on $\Lambda^2V$, and the manuscript's cusp fibre has the combinatorics of Mumford's degenerations of abelian surfaces;
\item \emph{thin groups}: the monodromy group is arithmetic, not thin.
\end{itemize}
The edges to the author's other workbenches, Yang--Mills, Erdős--Straus and Navier--Stokes, are listed with their status. The Erdős--Straus workbench uses exactly the dual action of Theorem~IV.

\subsection*{How this record was made}

The construction is that of the manuscript circulated by Levent Alpöge; the attribution file of the Yang--Mills workbench states that it was produced with Claude. The record is the work of KokunoYumeto, who designed the validate/repair/disprove programme and required proofs, calculations and counterexamples before any conclusion, with Codex for the verification and ChatGPT Pro for two independent reviews. Philip Engel's exposition is cited as \cite{EngelS6}.

This paper was prepared by Claude (Anthropic), model \texttt{claude-opus-5-5} (\emph{claude-ab} below), from an audit of the record on 26--27 September 2026. Independent Claude instances, each of which wrote none of the text, refereed the earlier version and the monodromy computations; they read for overstatement and for understatement and missed results, and several results were found in these passes, each credited at its statement. Every finding of another instance was verified by claude-ab before it was applied.

Every statement has one of four statuses:
\begin{itemize}
\item \emph{verified};
\item \emph{not verified here}, which implies nothing in either direction;
\item \emph{my assessment}, a view in the first person with its reasons, introduced by those words;
\item \emph{proved negative}, stated with its counterexample or derivation (Section~\ref{ssec:mononeg} collects those on the period system).
\end{itemize}
The status lines under the results make the first two precise. \textsf{Proved here} means that the proof is given in full in this paper; \textsf{Audited} means that the record's proof was read in full, every step checked, and independent checks run; \textsf{Reproduced} means that the finite content was recomputed and the general proof not re-derived. These are forms of \emph{verified}. \textsf{Located} means that the statement and its proof were found and are described, and \textsf{Stated} that the statement and the location of its proof are reported without the proof having been seen; both are forms of \emph{not verified here}. The record's own labels are \textsf{VERIFIED}, \textsf{OPEN} and \textsf{REFUTED}; they are reported as the record's, next to the status here.

The record's files were read as listed in Appendix~\ref{app:catalogue}: the correction note in full; the monograph's Part~I, ledger, Theorem~13.6 and §§13.11 and~14; the reconstructed manuscript's summary, Remarks~1.3 and~9.21, §2 (the period system) and §§10.1--10.3; the claim ledger; the guides and the key-advances reader. The rest was mapped from its tables of contents and summaries.

This is version~3. It keeps the theorems, proofs and citations of version~2 and corrects its framing, and it adds the integral monodromy of the period system (Section~\ref{sec:monodromy}) and its bridges. Version~2 and the audit notes cited as \note{NN\_} were withdrawn from the branch \texttt{claude/claude-ab-grind-20260925} on 27 September 2026 because of their framing; they remain, unchanged, in the branch's history at commit \texttt{a90d9e5b}, in \texttt{contrib/claude-ab/grind\_20260925/}. Appendix~\ref{app:changes} lists every change.

\subsection*{Organization}

Section~\ref{sec:problem} states the question and what the record claims, Section~\ref{sec:record} maps the record, Section~\ref{sec:candidate} outlines the construction, and Section~\ref{sec:ledger} reports the record's ledger. Section~\ref{sec:correction} audits the correction note and proves Theorem~I. Section~\ref{sec:later} treats the later supplements and the arithmetic of S6-5. Section~\ref{sec:monodromy} determines the monodromy group and proves Theorems~II--IV. Section~\ref{sec:overlaps} describes the bridges, and Section~\ref{sec:open} the open questions and what was not checked. The appendices give the catalogue of statements, the verification and the changes of version~3.


\tableofcontents
\clearpage

\section{The question, the published obstruction, and what the record claims}\label{sec:problem}

\subsection{The question}

Whether the six-sphere carries a complex structure (an integrable almost complex structure) has long been open. Campana, Demailly and Peternell studied a hypothetical complex structure $X$ on $S^6$ through its algebraic dimension $a(X)$, the transcendence degree of its field of meromorphic functions, first in 1998 \cite{CDP98} and then in 2020 \cite{CDP20}. In the 2020 paper they prove that a compact complex threefold homeomorphic to $S^6$ would have algebraic dimension $0$ (\cite{CDP20}, Corollary~2.3; Engel's exposition cites the same statement from the 1998 paper, \cite{CDP98}, Corollary~4.2). A complex structure on $S^6$ that fibres holomorphically over $\Pp^1$ has $a(X)\ge1$, and would therefore contradict that theorem.

\subsection{The claimed construction, the record, and an exposition}

\begin{itemize}
\item \emph{The source manuscript} \cite{AlpogeS6}, whose opening summary is titled \emph{The $(3,4,\infty)$ modular family of 2-tori, completed at its three special points, is a complex structure on $S^6$}, claims exactly such a threefold: $X\to\Pp^1$ with $a(X)=1$, $X$ diffeomorphic to $S^6$. The record reconstructs it as file \texttt{03\_} (108 pages). It was circulated by Levent Alpöge and produced with Claude, according to the attribution file of the Yang--Mills workbench.
\item \emph{The record} (24 August to 9 September 2026), by KokunoYumeto, audits that manuscript. It reconstructs the manuscript, repairs local defects, and concludes, with its own label \textsf{VERIFIED}, that the constructed threefold satisfies every hypothesis of the relevant theorem of \cite{CDP20} and contradicts its conclusion. It also isolates a step in the published proof that is not a formal consequence of the sequences displayed there. The source manuscript had already located the failure of that step at its own cusp fibre; the record's correction note adds an independent test family (Sections~\ref{sec:ledger} and~\ref{sec:correction}).
\item \emph{An exposition.} Philip Engel's \emph{Complex structures on $S^6$} (13 September 2026) \cite{EngelS6} presents the construction, which it credits to Claude under Alpöge's direction, with the stated aim of a self-contained proof. It notes that the construction contradicts the result that the algebraic dimension must be $0$, which it cites as \cite{CDP98}, Corollary~4.2, and it describes the construction as a one-parameter family of threefolds diffeomorphic to $S^6$.
\end{itemize}
The global construction and the counterexample claim are not verified here. This paper reports what the record claims, and proves or checks what can be established from the record's local, finite and algebraic statements.

\subsection{The record's own ledger and scope}

\begin{itemize}
\item At the freeze of version 1.0 its internal claim ledger (in the evidence archive \texttt{04\_}) had $42$ rows labelled \textsf{VERIFIED}, none \textsf{CONDITIONAL}, $2$ \textsf{OPEN} and $4$ \textsf{REFUTED}. The two open rows are:
\begin{itemize}
\item independent specialist replication of the cyclic invariant lattices, the logarithmic determinant divisor, the residue cancellation and an order-four scalar test;
\item specialist and community adjudication of the complete construction and counterexample.
\end{itemize}
The README describes them as replication and external adjudication, not missing internal premises.
\item The later supplements (5--9 September) use the construction's data as stated input, and state that they do not by themselves establish a complex structure on $S^6$ or disprove the published theorem.
\end{itemize}

\section{The record, and what each file contains}\label{sec:record}

\subsection{Four versions under one concept DOI}

The record 10.5281/zenodo.22678442 is the latest of four versions of the concept record 10.5281/zenodo.22235523. Each version kept the files of the one before and added new ones (read from the record's public metadata on 26 September 2026):

\begin{center}\small
\begin{tabularx}{\linewidth}{@{}>{\raggedright\arraybackslash}p{0.21\linewidth}>{\raggedright\arraybackslash}p{0.07\linewidth}>{\raggedright\arraybackslash}X@{}}
\toprule
Version (DOI 10.5281/zenodo.\dots) & Files added & What was added\\
\midrule
22235524, 1 Sept.\ (1.0) & 20 & The 154-page verification monograph (\fn{01\_}); the 128-page annotated manuscript (\fn{02\_}); the 108-page clean reconstruction of the source manuscript (\fn{03\_}); sources and evidence (\fn{04\_}, 20.8~MB); a mathematical provenance file and its receipt (\fn{05\_}, \fn{06\_}); conversation and workflow records, with receipts and validation reports (\fn{07\_}--\fn{13\_}); README, licence and rights files\\
22238779, 1 Sept.\ (1.1) & 2 & The eleven-page correction note on the singular-fibre step of \cite{CDP20}, with its source (\fn{00\_})\\
22346462, 5 Sept.\ (1.2) & 6 & A topology supplement: a reader's guide (\fn{14\_}), \emph{The higher torus wrap} on the cusp (\fn{15\_}, 50 pp.), \emph{Real $S^{12}$ and $S^{24}$ routes} (\fn{16\_}, 64 pp.), their sources and notes (\fn{17\_}--\fn{19\_})\\
22678442, 9 Sept.\ (consolidated) & 5 & The five-page key-advances reader of 6 September and its source (\fn{26\_}, \fn{27\_}); the complete frozen project of 6 September (\fn{28\_}, 174~MB, 1{,}079 files, including a 782-page cumulative workbench and a 138-page higher-rung paper); a guide and a manifest (\fn{29\_}, \fn{30\_})\\
\bottomrule
\end{tabularx}
\end{center}

The record's front file \texttt{README.md} was written for version 1.0 and describes it. The later guides are \fn{14\_s6\_topology\_update\_readers\_guide\_2026-09-05.pdf} and \fn{29\_S6\_FROZEN\_PROJECT\_GUIDE\_2026-09-09.md}.

\subsection{Where to start}

\begin{itemize}
\item \emph{For the claimed construction}: the clean reconstruction \fn{03\_} (108 pages). Its opening pages summarise the whole argument.
\item \emph{For the audit}: the monograph \fn{01\_}. Its Part~I (five pages) orients; its Part~II is the complete validate/repair/disprove paper, with its ledger and its dispositions (§§13.11 and 14).
\item \emph{For the published theorem}: the correction note \fn{00\_} (eleven pages), which is self-contained. Section~\ref{sec:correction} states and audits it.
\item \emph{For the later work}: the key-advances reader \fn{26\_}, and the guide \fn{14\_} to the topology supplement.
\end{itemize}

The files \fn{07\_}--\fn{13\_}, the record's conversation and workflow records with their receipts and validation reports, were not used here.

\section{The claimed construction, in outline}\label{sec:candidate}

This section follows the opening summary of the reconstructed source manuscript (file \fn{03\_}) and the monograph's orientation. Nothing in it was checked here, apart from the arithmetic noted.

\begin{enumerate}[label=\arabic*.]
\item \emph{Base.} The triangle group $\Delta=\Delta(3,4,\infty)\cong(\ZZ/3)*(\ZZ/4)$ acts on the upper half-plane. The quotient, with its cusp added, is $\Pp^1$ with orbifold points of orders $3$ and $4$ and one cusp.
\item \emph{A lattice representation.} On $V\cong\ZZ^4$ there are automorphisms $T_1$ of order $3$ and $T_2$ of order $4$ whose product has unipotent inverse $T_0$ with $(T_0-I)^2=0$; $\Delta$ acts on the dual lattice $\Lambda$ through them.
\item \emph{Periods and the family.} Period functions $\tau,\mu,\beta$ on the half-plane give an equivariant period map. The monograph writes the period matrix as
\[
\Pi=[Z\mid I_2],\qquad Z=\begin{pmatrix}6\mu&\tau\\ \beta&\mu\end{pmatrix},\qquad D=\operatorname{Im}\beta-\frac{6(\operatorname{Im}\mu)^2}{\operatorname{Im}\tau}<0,\qquad \det\nolimits_{\R}\Pi=\operatorname{Im}\tau\cdot D<0 .
\]
The quotient is a family of complex two-dimensional tori over the punctured orbifold curve. Its very general fibre is not algebraic, and its Néron--Severi group is generated by one indefinite class.
\item \emph{Three fillings.} At the two elliptic points, Kodaira's logarithmic transform: (disc $\times$ torus)$/(\ZZ/m)$, $m=3,4$, giving multiple fibres. At the cusp, Mumford's toric degeneration, whose central fibre is a degree-six del Pezzo surface glued to itself along opposite sides, with two triple points.
\item \emph{Gluing and the fundamental group.} The fillings are glued along collars, and the gluings can be changed by translations. The fundamental group and the homology depend on three integers $\ell_0,\ell_1,\ell_2$, with
\[
\pi_1(X)\cong\ZZ/|12\ell_0-4\ell_1-3\ell_2| .
\]
The record's candidate takes $(\ell_0,\ell_1,\ell_2)=(0,1,-1)$, for which $12\ell_0-4\ell_1-3\ell_2=-1$, so the fundamental group is trivial (source manuscript §7; monograph Part~I).
\item \emph{Conclusions claimed.} For this choice, the manuscript concludes:
\begin{itemize}
\item $X$ is simply connected with the integral homology of $S^6$;
\item $X$ is therefore diffeomorphic to $S^6$, as there are no exotic smooth structures on the six-sphere;
\item $c_3(X)=2$, localised at the cusp;
\item $a(X)=1$, with the fibration as algebraic reduction.
\end{itemize}
\end{enumerate}

The period data of item~3 are also what the owner's Yang--Mills and Erdős--Straus workbenches take from this construction (Section~\ref{sec:overlaps}). In the Yang--Mills record the quantity $D_{\mathrm{YM}}=L_Tq_T+6m_T^2$, with $L_T=\operatorname{Im}\tau$, $q_T=-\operatorname{Im}\beta$ and $m_T=\operatorname{Im}\mu$, equals $-\operatorname{Im}\tau\cdot D=-\det_{\R}\Pi>0$. So the two records use the same number, with opposite sign conventions.

\section{The record's ledger and its central claims}\label{sec:ledger}

The record labels each claim with its own status: \textsf{VERIFIED} for a complete proof, \textsf{OPEN} when neither a proof nor a counterexample is recorded, and \textsf{REFUTED} for an exact assertion with a recorded contradiction or counterexample. Two lists carry them, and this paper cites both:
\begin{itemize}
\item the \emph{claim ledger} in the evidence archive \texttt{04\_}, with 48 rows: 42 \textsf{VERIFIED}, 2 \textsf{OPEN}, 4 \textsf{REFUTED}. The README quotes its counts.
\item the monograph's \emph{status list}, the 21 numbered items of its Part~I. Item numbers below refer to this list.
\end{itemize}

\subsection{The central claims (record: \textsf{VERIFIED})}

\begin{itemize}
\item \emph{The construction.} According to the monograph, its sections 6--7 prove the period-family descent, the toric and finite fillings, the global gluing, and the integral nearby-cycle and Leray comparisons. With them the constructed threefold $X$ satisfies $H^1(X;\ZZ)=H^2(X;\ZZ)=0$, $a(X)=1$ and $c_3(X)=2$, and $f:X\to\Pp^1$ is holomorphic and an algebraic reduction.
\item \emph{A counterexample to the published theorem.} $X$ satisfies every stated hypothesis of Theorem~2.2 of \cite{CDP20} and contradicts its conclusion (monograph §13.11, item~5). The monograph says that this uses the complete construction and topology proofs, not only the two defects below.
\item \emph{A corrected criterion.} This is the part of the published argument that survives. If some line bundle $L_0$ has $H^2(X,T_X\otimes L_0)=0$, then generic vanishing holds and $c_3(X)\le0$ (monograph Theorem 13.28; the correction note's Theorem 5.2). For the constructed $X$ this hypothesis fails for every $L_0$, since $c_3(X)=2$.
\item \emph{Further invariants}: the Hodge diamond and the Frölicher sequence, the Bott--Chern and Aeppli dimensions, and global representatives with the Bott--Chern product structure (status list, items 17, 19 and 20).
\end{itemize}
\status{\textsf{Located}; not verified here. The record lists their independent replication and adjudication as its two open rows (Section~\ref{ssec:openrows}).}

\subsection{The four rows labelled \textsf{REFUTED}}

Each row is given with its claim-ledger identifier.

\begin{enumerate}[label=\arabic*.]
\item \emph{The singular-fibre inference of \cite{CDP20}} (p.~692 of the journal version; \texttt{S6-CDP-P692}; status list, item~7).
\begin{itemize}
\item The inference: for an irreducible reduced fibre, the published proof reduces a vanishing to the vanishing of $H^0(S,\widetilde\Omega^1_S\otimes L|_S)$.
\item The record says that this is not a formal consequence of the two cotangent sequences, and gives a test family.
\item The source manuscript had already located the failure at its own cusp fibre (its Lemma~10.3 and Theorem~10.5), and had noted that the step is valid on a normal fibre (its Remark~10.4).
\end{itemize}
\status{The test family and the corrected criterion are \textsf{audited} here. The exact size of the failure at fibres with normal crossings, and the exact condition for the torsion that causes it, are \textsf{proved here} (Section~\ref{sec:correction}). The source manuscript quotes the published sentence verbatim (its §10.1). The paper itself was not re-read here.}
\item \emph{A homology formula imported from \cite{CDP98}} (\texttt{S6-CDP98-MONO}; status list, item~12). The direct-product formula of the 1998 paper (its Lemma~3.2, in the record's reading) omits the coinvariants of the monodromy. This leaves a torus branch of the 2020 argument without its printed rank argument. \status{\textsf{Located}.}
\item \emph{A remark of the source manuscript on relative differentials} (\texttt{S6-HODGE-CUSP-TORSION}; status list, item~18).
\begin{itemize}
\item The manuscript's Remark~1.3 speaks of the sheaves $\Omega^p_{X/B}$, and its Remark~9.21 of $\Omega^1_{X/B}$, as having torsion along all three singular fibres.
\item The record shows that at the cusp $\Omega^1_{X/B}$ is torsion-free but not locally free, and that its torsion occurs only at the two multiple fibres.
\end{itemize}
\status{\textsf{Proved here} for the local models the record gives (Proposition~\ref{prop:torsion}). The refutation is correct for $p=1$. For $p=2$ the same local models do have torsion along the double curves of the cusp fibre, so Remark~1.3's statement holds for $\Omega^2_{X/B}$.}
\item \emph{A remark of the source manuscript on the gluing integers} (\texttt{S6-LOCAL-728}; monograph §14, item~2). The remark (its Remark~7.28) is stated for a fixed triple. The triple $(\ell_0,\ell_1,\ell_2)=(100,1,-1)$ gives $12\ell_0-4\ell_1-3\ell_2=1199\neq\pm1$, so the statement must quantify existentially over $\ell_0$ for fixed admissible $(\ell_1,\ell_2)$. \status{\textsf{Reproduced}: the arithmetic.}
\end{enumerate}

\subsection{The two rows labelled \textsf{OPEN}}\label{ssec:openrows}

\begin{itemize}
\item \texttt{S6-OPEN-HODGE-REPLICATION}: independent specialist replication of the cyclic invariant lattices, the logarithmic determinant divisor, the residue cancellation and an order-four scalar test. The monograph's status list has this as its single open item (item~21), together with the three cotangent direct images and the values $h^{1,2}=1$, $h^{1,1}=2$.
\item \texttt{S6-OPEN-EXTERNAL-ADJUDICATION}: specialist and community adjudication of the complete construction and counterexample.
\end{itemize}
The monograph states that these rows record replication and adjudication that are outstanding, and that they are not unresolved premises of its proofs.

\subsection{Confirmed local repairs}

Monograph §14 lists seven items, labelled separately from the theorem. Five are labelled \textsf{VERIFIED}:
\begin{itemize}
\item a closed-image argument;
\item an appendix sequence whose arrows are cohomological although it is labelled as homology;
\item a fundamental-group identification;
\item a build defect of the record's own assembler;
\item a correction to a finite-filling computation of a canonical section.
\end{itemize}
The other two are the \textsf{REFUTED} remarks treated above. The five are \textsf{located}.

\section{The correction note on the singular-fibre step, audited}\label{sec:correction}

The correction note (file \fn{00\_}, eleven pages, dated 1 September 2026, added in version 1.1) isolates one step in the singular-fibre part of the proof in \cite{CDP20} and records what survives without it. Its scope statement is narrow: it does not infer that Theorem~2.2 of \cite{CDP20} is false, and it credits the source manuscript for the conductor-supported differential on which it builds (the manuscript's §10.2, Lemma~10.3). The source manuscript had already located the failure of this step at its own cusp fibre (its remark on \cite{CDP20} after the main theorem, and its Theorem~10.5). In its Remark~10.4 it had also recorded that the step is valid when the fibre is normal.

This section states the note's results and reports the audit. It adds four things, each proved here:
\begin{itemize}
\item the exact condition for the torsion that makes the step fail: it is present exactly at the non-normal points (Proposition~\ref{prop:normal});
\item a description of the kernel sheaf on which the step depends. At a reduced fibre with normal crossings, the group that the step asserts to vanish is the space of sections of the line bundle pulled back to the normalization (Proposition~\ref{prop:kernel}). This was found in the first referee pass.
\item the exact size of the failure in the note's test family (Propositions~\ref{prop:sharp} and~\ref{prop:kernel});
\item the general form of the note's local computation of relative differentials, including second exterior powers (Proposition~\ref{prop:torsion}).
\end{itemize}
\status{\textsf{Audited}. Every proof in the note was read in full and every step checked, and the finite and symbolic content was recomputed (\fn{s6r\_checks.py}, Appendix~\ref{app:verification}). No mathematical error was found.}

\subsection{The filtration and the exact criterion}

Let $X$ be a complex manifold and $S\subset X$ a reduced effective Cartier divisor, with smooth locus $j:S_{\mathrm{reg}}\hookrightarrow S$ and conormal sheaf $N^*=N^*_{S/X}$. Put
\[
T_S:=\ker\bigl(\Omega^1_S\to j_*\Omega^1_{S_{\mathrm{reg}}}\bigr),\qquad \widetilde\Omega^1_S:=\Omega^1_S/T_S,\qquad K_S:=\ker\bigl(\Omega^1_X|_S\to\widetilde\Omega^1_S\bigr).
\]
By the note's Lemma~1.1, $T_S$ is the torsion subsheaf of $\Omega^1_S$.

\begin{proposition}[the note's Lemma 2.1 and Propositions 2.2--2.3]\label{prop:filtration}
\leavevmode
\begin{enumerate}[label=\textup{(\alph*)}]
\item The conormal sequence $0\to N^*\to\Omega^1_X|_S\to\Omega^1_S\to0$ is exact on the left.
\item There are exact sequences $0\to N^*\to K_S\to T_S\to0$ and $0\to K_S\to\Omega^1_X|_S\to\widetilde\Omega^1_S\to0$.
\item Let $A$ be a vector bundle on $S$ with $H^0(S,\widetilde\Omega^1_S\otimes A)=0$ and $H^0(S,N^*\otimes A)=0$. Then
\[
H^0(S,\Omega^1_X|_S\otimes A)\;\cong\;\ker\Bigl(\delta_A:H^0(S,T_S\otimes A)\to H^1(S,N^*\otimes A)\Bigr),
\]
where $\delta_A$ is the connecting map of the first sequence in \textup{(b)} tensored with $A$.
\end{enumerate}
\end{proposition}

\begin{proof}
(a) Locally $S=\{g=0\}$ with $g$ reduced in a regular local ring, and the map is $\bar h\mapsto h\,dg$. The differential $dg$ is nonzero at the general point of every component of $S$, since each component is generically smooth and $g$ generates its ideal there. So $h\,dg\in g\,\Omega^1$ forces $h$ to vanish on every component, and $h\in(g)$ because $(g)$ is radical.
(b) $K_S$ is the preimage of $T_S$ under the surjection $\Omega^1_X|_S\to\Omega^1_S$, whose kernel is $N^*$ by (a).
(c) Tensor both sequences of (b) with $A$ and take cohomology. The second gives $H^0(\Omega^1_X|_S\otimes A)\cong H^0(K_S\otimes A)$, and the first gives $H^0(K_S\otimes A)\cong\ker\delta_A$.
\end{proof}

\status{\textsf{Audited}. The note states (c) for a line bundle on a threefold; the proof uses neither restriction. Reducedness is needed in (a): for $g=g_0^2$, the class of $g_0\cdot g$ in $N^*$ is nonzero, and it maps to $g_0\,dg=2g\,dg_0$, which is zero in $\Omega^1_X|_S$.}

Thus the implication at issue,
\[
H^0(S,N^*\otimes A)=0\ \text{ and }\ H^0(S,\widetilde\Omega^1_S\otimes A)=0\quad\Longrightarrow\quad H^0(S,\Omega^1_X|_S\otimes A)=0,
\]
holds whenever $T_S=0$, since then $K_S=N^*$. In general it holds exactly when $\delta_A$ is injective (under its two hypotheses).

\subsection{Where the torsion lives}

\begin{lemma}[modules with one relation]\label{lem:onerel}
Let $R$ be a Noetherian ring with total ring of fractions $Q$, and let $v\in R^n$ be a vector whose entries generate an ideal $J$ containing a nonzerodivisor. Then $\varphi\mapsto[\varphi v]$ induces an isomorphism
\[
(R:_QJ)/R\;\xrightarrow{\ \sim\ }\;\operatorname{Tors}\bigl(R^n/Rv\bigr),\qquad\text{and}\qquad (R:_QJ)/R\;\cong\;\operatorname{Ext}^1_R(R/J,R).
\]
In particular, when $J\neq R$, the module $R^n/Rv$ is torsion-free if and only if $J$ has grade at least $2$; when $J=R$ it is free.
\end{lemma}

\begin{proof}
For $\varphi\in(R:_QJ)$ the vector $\varphi v$ lies in $R^n$. If $f$ is a nonzerodivisor with $f\varphi\in R$, then $f[\varphi v]=[(f\varphi)v]=0$, so the class is torsion. It is zero iff $\varphi v=cv$ for some $c\in R$, that is, iff $(\varphi-c)J=0$ in $Q$; since $J$ contains a unit of $Q$, this means $\varphi=c\in R$. Conversely, if $f[a]=0$ for a nonzerodivisor $f$, then $fa=hv$ with $h\in R$; then $\varphi=h/f$ satisfies $\varphi v=a\in R^n$, so $\varphi\in(R:_QJ)$ and $[a]$ is its image.

For the second isomorphism, apply $\operatorname{Hom}_R(-,R)$ to $0\to J\to R\to R/J\to0$. Here $\operatorname{Hom}(R/J,R)=0$, since $J$ contains a nonzerodivisor, and $\operatorname{Ext}^1(R,R)=0$. Every homomorphism $\psi:J\to R$ is multiplication by $\psi(f_0)/f_0\in Q$ for any nonzerodivisor $f_0\in J$. So $\operatorname{Hom}(J,R)=(R:_QJ)$, and the sequence gives $(R:_QJ)/R\cong\operatorname{Ext}^1(R/J,R)$. By Rees's theorem \cite[Thm.~16.6]{Mat86}, the grade of $J$ is the least $i$ with $\operatorname{Ext}^i(R/J,R)\neq0$. The grade is at least $1$, and it is at least $2$ iff $\operatorname{Ext}^1$ vanishes.
\end{proof}

\begin{proposition}[torsion of the differentials of a hypersurface]\label{prop:normal}
Let $S=\{g=0\}$ be a reduced hypersurface in a complex manifold of dimension $n\ge2$. Let $x\in S$, let $R=\mathcal O_{S,x}$, and let $J\subset R$ be the ideal generated by the restrictions of $\partial g/\partial z_1,\dots,\partial g/\partial z_n$. Then
\[
T_{S,x}\;\cong\;\operatorname{Ext}^1_R(R/J,R),
\]
and $T_{S,x}=0$ iff $\operatorname{ht}J\ge2$. Equivalently, $T_{S,x}=0$ iff the singular locus of $S$ has codimension at least $2$ in $S$ at $x$, that is, iff $S$ is normal at $x$. When a component of $\operatorname{Sing}S$ through $x$ has codimension one, $T_{S,x}\neq0$. Globally, let $\mathcal J\subset\mathcal O_S$ be the Jacobian ideal sheaf, generated locally by the restrictions of the $\partial g/\partial z_i$. Then $\varphi\otimes[g]\mapsto[\varphi\,dg]$ is a canonical isomorphism
\[
\mathcal{E}xt^1_{\mathcal O_S}(\mathcal O_S/\mathcal J,\mathcal O_S)\otimes N^*_{S/X}\;\xrightarrow{\ \sim\ }\;T_S .
\]
\end{proposition}

\begin{proof}
The differentials have the presentation $\Omega^1_{S,x}=R^n/Rv$ with $v=(\partial g/\partial z_i|_S)_i$.

\emph{$J$ contains a nonzerodivisor.} The ring $R$ is reduced, so its zerodivisors form the union of its minimal primes. No minimal prime contains $J$, because $dg\neq0$ at the general point of each component. Prime avoidance then gives an element of $J$ outside all minimal primes.

\emph{The torsion.} Lemma~\ref{lem:onerel} gives $T_{S,x}\cong\operatorname{Ext}^1(R/J,R)$, which vanishes iff $\operatorname{grade}J\ge2$.

\emph{Grade, height and codimension.} $R$ is a regular local ring modulo a nonzerodivisor, hence Cohen--Macaulay. So $\operatorname{grade}J=\operatorname{ht}J=\dim R-\dim R/J$ \cite[Thm.~17.4]{Mat86}. A point $y\in S$ is smooth iff $dg(y)\neq0$: if $S$ is smooth at $y$, its ideal is generated by one coordinate function $w$, so $g$ is a unit times $w$. Hence the zero set of $J$ is $\operatorname{Sing}S$, and $\dim R/J$ is the dimension of the germ of $\operatorname{Sing}S$ at $x$ \cite{GR84}. Since $R$ is Cohen--Macaulay, it satisfies Serre's condition $S_2$. So, by Serre's criterion \cite[Thm.~23.8]{Mat86}, it is normal iff it is regular in codimension one, that is, iff the singular locus has codimension at least $2$ (for normal complex spaces see also \cite{GR84}).

\emph{The global form.} Replace $g$ by $ug$ with $u$ a unit. On $S$, $\partial(ug)/\partial z_i=u\,\partial g/\partial z_i$, so $\mathcal J$ does not change. The form $\varphi\,dg$, for $\varphi\in(\mathcal O_S:_Q\mathcal J)$, is holomorphic and does not depend on coordinates, and $(\varphi/u)\,d(ug)=\varphi\,dg$ on $S$. So the local isomorphisms of Lemma~\ref{lem:onerel}, $\varphi\mapsto[\varphi\,dg]$, glue once the generator $[g]$ of $N^*_{S/X}$ is carried along. The identification $(R:_QJ)/R\cong\operatorname{Ext}^1_R(R/J,R)$ of Lemma~\ref{lem:onerel} is canonical.
\end{proof}

\begin{corollary}\label{cor:normal}
The implication at issue holds on every normal reduced fibre (for a surface fibre: one whose singular points are isolated). On a non-normal reduced fibre, $T_S\neq0$, and the implication holds exactly when $\delta_A$ is injective.
\end{corollary}

The first statement is the source manuscript's Remark~10.4, proved there by the second Riemann extension theorem. Proposition~\ref{prop:normal} adds the converse: on a reduced hypersurface the torsion of $\Omega^1_S$ is present at exactly those points where the singular locus has a component of codimension one. The fibre of the candidate threefold to which the step would be applied is the cusp fibre, which the record describes as a degree-six del Pezzo surface glued to itself along curves. The corollary therefore does not cover it, as the record says. On any reduced fibre $N^*_{S/X}\cong\mathcal O_S$ (on a multiple fibre it is a nontrivial torsion bundle), so the map on which the step depends is
\[
\delta_A:\ H^0\bigl(S,\ \mathcal{E}xt^1_{\mathcal O_S}(\mathcal O_S/\mathcal J,\mathcal O_S)\otimes A\bigr)\longrightarrow H^1(S,A),
\]
a map between groups defined by the fibre and $A$ alone. For a fibre with normal crossings, such as the candidate's cusp fibre, $\mathcal J$ is the reduced ideal of the double curve.

For $p=1$, Proposition~\ref{prop:normal} is a case of theorems of Vetter \cite{Vet70} and Lebelt \cite{Leb74}, summarized in \cite{MV19}. Let a reduced complete intersection have singular locus of codimension $k$. Then $\Omega^p$ is torsion-free for $p<k$ and has torsion for $p=k$ (Vetter), and conversely torsion-freeness forces $p<k$ (Lebelt). The proof above is elementary and gives the torsion as an $\operatorname{Ext}$ sheaf.

\begin{proposition}[the kernel sheaf]\label{prop:kernel}
Let $S=\{g=0\}$ be a reduced hypersurface in a complex manifold $X$, with Jacobian ideal sheaf $\mathcal J$.
\begin{enumerate}[label=\textup{(\alph*)}]
\item The map $\varphi\otimes[g]\mapsto\varphi\,dg$ is an isomorphism $\mathcal{H}om_{\mathcal O_S}(\mathcal J,\mathcal O_S)\otimes N^*\to K_S$, which restricts to $\mathcal O_S\otimes N^*\cong N^*$. So the first sequence of Proposition~\ref{prop:filtration}(b) is the sequence $0\to\mathcal O_S\to\mathcal{H}om(\mathcal J,\mathcal O_S)\to\mathcal{E}xt^1(\mathcal O_S/\mathcal J,\mathcal O_S)\to0$, tensored with $N^*$.
\item Suppose $S$ has normal crossings: locally $g$ is a unit times $x$, $xy$ or $xyz$ in suitable coordinates. Let $\eta:\tilde S\to S$ be the normalization. Then $\mathcal{H}om(\mathcal J,\mathcal O_S)=\eta_*\mathcal O_{\tilde S}$. Hence, if $H^0(S,\widetilde\Omega^1_S\otimes A)=0$ for a line bundle $A$ on $S$, then
\[
H^0(S,\Omega^1_X|_S\otimes A)\;\cong\;H^0\bigl(\tilde S,\eta^*(N^*\otimes A)\bigr).
\]
\item On a reduced fibre with normal crossings and with $H^0(S,\widetilde\Omega^1_S\otimes A)=0$, the conclusion $H^0(S,\Omega^1_X|_S\otimes A)=0$ of the step holds if and only if $H^0(\tilde S,\eta^*A)=0$.
\item In the note's test family, $H^0(S,\Omega^1_{\mathcal X}|_S\otimes A)\cong\C$ for every $\lambda\in\C^*$. For $\lambda=1$ it is spanned by $dt$.
\end{enumerate}
\end{proposition}

\begin{proof}
(a) Work at $x$ with $R$ and $v$ as in Proposition~\ref{prop:normal}, and identify $\Omega^1_X|_S$ at $x$ with $R^n$ through the $dz_i$.
\begin{itemize}
\item A vector $a$ lies in $K_{S,x}$ iff $[a]$ is torsion.
\item By Lemma~\ref{lem:onerel} this holds iff $[a]=[\varphi v]$ for some $\varphi\in(R:_QJ)$. Since $R\subseteq(R:_QJ)$, this holds iff $a\in(R:_QJ)\,v$.
\item The map $\varphi\mapsto\varphi v$ is injective, since $J$ contains a unit of $Q$.
\end{itemize}
So $K_{S,x}=(R:_QJ)\,dg\cong(R:_QJ)=\operatorname{Hom}_R(J,R)$, and $R\,dg$ is $N^*$. The choices do not matter, as in Proposition~\ref{prop:normal}.

(b) Choose coordinates in which $g=x_1\cdots x_k$ with $k\le3$, the unit being absorbed into $x_1$. The branches are $S_i=\{x_i=0\}$, and $\tilde R=\prod_i\C\{\hat x_i\}$ is the ring of tuples of functions on the branches.
\begin{itemize}
\item \emph{$R$ inside $\tilde R$.} $R$ consists of the tuples that agree on the pairwise intersections $S_i\cap S_l$. For $k=2$ this is the note's Lemma~3.1, for $k=3$ the manuscript's Lemma~10.2.
\item \emph{The generators of $J$.} The generator $\prod_{j\ne i}x_j$ of $J$ vanishes on every branch but $S_i$. So, for $\varphi=(\varphi_l)\in Q$, the product $\varphi\prod_{j\ne i}x_j$ lies in $R$ iff $\varphi_i\prod_{j\ne i}x_j$ is holomorphic on $S_i$ and vanishes on each $S_i\cap S_l$.
\item \emph{The quotient.} The $x_j$ with $j\ne i$ are distinct primes of $\C\{\hat x_i\}$, so this holds iff $\varphi_i$ is holomorphic. Hence $(R:_QJ)=\tilde R$. For $k=1$ both sides are $R$.
\end{itemize}
The projection formula for the finite map $\eta$ gives $H^0(S,K_S\otimes A)=H^0(\tilde S,\eta^*(N^*\otimes A))$. Proposition~\ref{prop:filtration} identifies this group with $H^0(S,\Omega^1_X|_S\otimes A)$ when $H^0(S,\widetilde\Omega^1_S\otimes A)=0$.

(c) On a reduced fibre $N^*\cong\mathcal O_S$.

(d) In the test family, $\eta^*A$ is trivial on the connected surface $\Pp^1\times\Pp^1$. Also $H^0(S,\widetilde\Omega^1_S\otimes A)\hookrightarrow H^0(\Pp^1\times\Pp^1,\Omega^1)=0$ for every $\lambda$, by the proof of Theorem~\ref{thm:testfamily}. For $\lambda=1$ the bundle $A$ is trivial, and $dt$ is a nonzero section.
\end{proof}

For the candidate's cusp fibre $W$, whose normalization is a degree-six del Pezzo surface, the record states two things for every $L$:
\begin{itemize}
\item $\eta^*A\cong\mathcal O$ (the manuscript's Theorem~10.5, Step~1);
\item $H^0(W,\widetilde\Omega^1_W\otimes A)=0$ (the manuscript's remark on \cite{CDP20} after its main theorem).
\end{itemize}
If both hold, Proposition~\ref{prop:kernel}(b) gives $H^0(W,\Omega^1_X|_W\otimes A)\cong\C$ for every $L$: the failure there would have size exactly one.

\status{Lemma~\ref{lem:onerel}, Propositions~\ref{prop:normal} and~\ref{prop:kernel}, and Corollary~\ref{cor:normal}: \textsf{proved here}. Proposition~\ref{prop:normal} is classical in general form. Proposition~\ref{prop:kernel} was found in the first referee pass. The statement about $W$ is conditional on the record.}

\subsection{The test family, audited and sharpened}

Let $\Delta$ be the disc $|t|<4/27$, let
\[
\mathcal C=\{([x:y:z],t)\in\Pp^2\times\Delta:\ y^2z=x^2(x+z)+tz^3\},
\]
put $\mathcal X=\mathcal C\times\Pp^1$ with $g=t$, and let $S=C_0\times\Pp^1$ be the central fibre. Here $C_0$ is the nodal cubic, with node $n=[0:0:1]$. The double curve of $S$ is $D=\{n\}\times\Pp^1$, and the normalization is $\eta=\nu\times\mathrm{id}:\Pp^1\times\Pp^1\to S$. For $\lambda\in\C^*$, let $A_\lambda$ be the line bundle on $C_0$ obtained by gluing the fibres of the trivial bundle over the two preimages $p=[1:1]$ and $q=[-1:1]$ of the node by $\lambda$, and put $A=\mathrm{pr}_1^*A_\lambda$.

\begin{theorem}[the note's Theorem 4.4]\label{thm:testfamily}
If $\lambda$ is not a root of unity, then $H^0(S,N^*\otimes A)=0$ and $H^0(S,\widetilde\Omega^1_S\otimes A)=0$, but $H^0(S,\Omega^1_{\mathcal X}|_S\otimes A)\neq0$.
\end{theorem}

Every step of the note's proof was checked. The recomputed parts are:
\begin{itemize}
\item \emph{The total space and the fibres.} $\mathcal C$ is smooth: on $z\ne0$, $\partial F/\partial t=-z^3$; on $z=0$, $\partial F/\partial z=y^2\ne0$. The discriminant of $x^3+x^2+t$ is $-t(27t+4)$, so the only singular fibre over $\Delta$ is $t=0$.
\item \emph{The node and its normalization.} $\nu[u:v]=[(u^2-v^2)v:u(u^2-v^2):v^3]$ parametrizes $C_0$ and sends $p$ and $q$ to the node. At the node, $t=UV$ with $U=y-x\sqrt{1+x}$ and $V=y+x\sqrt{1+x}$, whose Jacobian determinant at the origin is $-2$.
\item \emph{The gluing bundle.} $H^0(C_0,A_\lambda)=0$ for $\lambda\ne1$. With the sections $O(t)=[0:1:0]$ and $P(t)=[3:\sqrt{36+t}:1]$, the bundle $\mathcal O_{\mathcal C}(P-O)$ restricts to $A_{1/3}$ (or $A_3$, by convention): the multiplier is $f(1)/f(-1)=1/3$ for $f(r)=r-2$, since $P(0)$ has $r=y/x=2$. So a non-torsion $A$ that extends to $\mathcal X$ exists (the note's Lemma~4.3).
\item \emph{The torsion at a crossing.} Take $R=\C\{u,v,w\}/(uv)$ and $\tau=v\,du=-u\,dv$. The kernel of $\Omega^1_R\to$ (the differentials of the two branches) is $R\tau$, with $\tau\neq0$ and $(u+v)\tau=0$. Hence $\widetilde\Omega^1_S\hookrightarrow\eta_*\Omega^1_{\Pp^1\times\Pp^1}$, and $H^0(\widetilde\Omega^1_S\otimes A)\hookrightarrow H^0(\Pp^1\times\Pp^1,\Omega^1)=0$.
\item \emph{The conductor differential.} On each branch, the pullback of $dg=v\,du+u\,dv$ vanishes along the preimage of $D$. So the conductor sequence (the note's Lemma~3.1, checked in the local model) gives a unique $s\neq0$ in $H^0(S,\Omega^1_{\mathcal X}|_S\otimes A)$ with $\eta^*s=\eta^*dg\otimes e$ and $s|_D=0$ (Proposition~3.2).
\end{itemize}

The monograph's version of this family (its Theorem~13.6, version 1.0) states that no extension of $A$ to the total space is claimed. The note, added in version 1.1, supplies one; the multiplier $1/3$ is reproduced above.

\begin{proposition}[the exact size of the failure]\label{prop:sharp}
In the test family, let $\lambda\neq1$. Then:
\begin{enumerate}[label=\textup{(\alph*)}]
\item $T_S\cong\mathcal O_D$, generated by $\tau=v\,du$;
\item $H^0(S,A)=H^1(S,A)=0$;
\item $\delta_A=0$, and $H^0(S,\Omega^1_{\mathcal X}|_S\otimes A)\cong H^0(S,T_S\otimes A)\cong\C$, spanned by the conductor differential $s$;
\item the same holds for the curve alone: $H^0(C_0,\Omega^1_{\mathcal C}|_{C_0}\otimes A_\lambda)\cong\C$.
\end{enumerate}
\end{proposition}

\begin{proof}
(a) At a point of $D$ take coordinates $(u,v,w)$ with $w$ a coordinate on $\Pp^1$, so that $\mathcal O_{S}=R=\C\{u,v,w\}/(uv)$. The note shows that the torsion there is $R\tau$. To find its annihilator, suppose $c\,v\,du=h(v\,du+u\,dv)$ in $R\,du\oplus R\,dv\oplus R\,dw$. Then $hu=0$ and $(c-h)v=0$. Since $\operatorname{Ann}(u)=(v)$ and $\operatorname{Ann}(v)=(u)$, this gives $c\in(u,v)$; conversely $u\tau=uv\,du=0$ and $v\tau=-uv\,dv=0$. So $R\tau\cong R/(u,v)=\mathcal O_{D}$. The functions $u,v$ come from the surface $\mathcal C$ and do not depend on the factor $\Pp^1$, so $\tau$ is defined on a whole neighbourhood of $D$ in $S$, and $T_S=\mathcal O_S\tau\cong\mathcal O_D$.

(b) Tensor the conductor sequence with $A$:
\[
0\to A\to\eta_*\eta^*A\oplus A|_D\to\eta_*\bigl(\eta^*A|_{D_1\sqcup D_2}\bigr)\to0,
\]
where $D_1=\{p\}\times\Pp^1$ and $D_2=\{q\}\times\Pp^1$. Here $\eta^*A\cong\mathcal O$ with a frame $e$, and $A|_D\cong\mathcal O_D$ with a frame $\sigma$. The ratios $e|_{D_i}/\eta^*\sigma$ are nowhere-vanishing holomorphic functions on $\Pp^1$, hence constants $\alpha,\beta$, and $\alpha/\beta=\lambda^{\pm1}$ by the construction of $A_\lambda$. Moreover $H^0(\Pp^1\times\Pp^1,\mathcal O)=H^0(\Pp^1,\mathcal O)=\C$ and $H^1(\Pp^1\times\Pp^1,\mathcal O)=H^1(\Pp^1,\mathcal O)=0$ \cite[Thm.~III.5.1]{Har77}, with GAGA \cite{Ser56} for the analytic groups; for the quadric $\Pp^1\times\Pp^1\subset\Pp^3$ use $0\to\mathcal O_{\Pp^3}(-2)\to\mathcal O_{\Pp^3}\to\mathcal O_Q\to0$. The long exact sequence therefore reads
\[
0\to H^0(S,A)\to\C\oplus\C\xrightarrow{\ \mu\ }\C\oplus\C\to H^1(S,A)\to0,\qquad \mu(c,k)=(\alpha c-k,\ \beta c-k).
\]
Since $\det\mu=\beta-\alpha\neq0$ for $\lambda\neq1$, both groups vanish.

(c) Since $S$ is a fibre, $N^*\cong\mathcal O_S$. So $\delta_A$ maps into $H^1(S,A)=0$, and $\delta_A=0$. The two endpoint groups vanish: the first by (b), the second as in the note. Proposition~\ref{prop:filtration}(c) then gives $H^0(\Omega^1_{\mathcal X}|_S\otimes A)\cong H^0(S,T_S\otimes A)=H^0(D,A|_D)=\C$, which contains $s\neq0$.

(d) The same argument on $C_0$ gives $T_{C_0}=\C_n\tau$ and $H^i(C_0,A_\lambda)=0$ for $i=0,1$. Hence the dimension is $1$.

An independent count on the normalization agrees. A section of $E=\Omega^1_{\mathcal C}|_{C_0}\otimes A_\lambda$ is a section $\tilde\sigma$ of $\nu^*E\cong\nu^*\Omega^1_{\mathcal C}$ whose values at $p$ and $q$ agree up to the factor $\lambda$ in $T^*_n\mathcal C$. There is an exact sequence $0\to N^*_\nu\to\nu^*\Omega^1_{\mathcal C}\to\Omega^1_{\Pp^1}\to0$, since $\nu$ is an immersion. The conormal bundle $N^*_\nu$ has degree $0-(-2)=2$, because $\deg\nu^*K_{\mathcal C}=K_{\mathcal C}\cdot C_0=\deg\omega_{C_0}-C_0^2=0$ by adjunction. Since $H^0(\Pp^1,\Omega^1)=0$, every global $\tilde\sigma$ is a section of $N^*_\nu$, so it takes values in the conormal lines of the two branches. These lines are distinct at the node, since $dU$ and $dV$ are independent. The gluing condition therefore forces $\tilde\sigma(p)=\tilde\sigma(q)=0$, and $\tilde\sigma\in H^0(N^*_\nu(-p-q))\cong H^0(\Pp^1,\mathcal O)=\C$.
\end{proof}

\begin{remark}\label{rem:sharp}
So in the test family $\delta_A$ is not merely non-injective: its target is zero, and every torsion section lifts. Only $\lambda\neq1$ is needed. The condition that $\lambda$ is not a root of unity serves only to meet the non-torsion proviso of the published step. The factor $\Pp^1$ serves only to reach dimension three: the phenomenon is already present on the nodal fibre of a family of curves.
\end{remark}

Proposition~\ref{prop:kernel}(d) gives the same dimension by a second route, and for every $\lambda\in\C^*$, including $\lambda=1$.

\status{Theorem~\ref{thm:testfamily}: \textsf{audited}. Proposition~\ref{prop:sharp}: \textsf{proved here}; the record states only the nonvanishing.}

\subsection{What survives: the derived global criterion}

\begin{theorem}[the note's Theorem 5.2; the monograph's Theorem 13.28]\label{thm:derived}
Let $X$ be a connected compact complex threefold with $H^1(X;\ZZ)=H^2(X;\ZZ)=0$ and $a(X)=1$, whose algebraic reduction is a holomorphic map onto a smooth compact curve. If some $L_0\in\operatorname{Pic}(X)$ has $H^2(X,T_X\otimes L_0)=0$, then $H^i(X,T_X\otimes M)=0$ for $i=0,2,3$ and all $M$ in a dense open subset of $\operatorname{Pic}(X)$, and $\langle c_3(T_X),[X]\rangle\le0$. By the note's Corollary~5.3 it suffices that $R^2f_*(T_X\otimes L)=0$ for one $L$.
\end{theorem}

The proof in the note was checked step by step:
\begin{itemize}
\item $\operatorname{Pic}(X)\cong H^1(X,\mathcal O_X)$ by the exponential sequence.
\item Sections vanishing to all orders along a fibre divisor vanish identically (Krull's intersection theorem), which gives $H^0(T_X(-mD))=0$ and $H^0(\Omega^1_X\otimes K_X(-\ell D))=0$ for large $m,\ell$.
\item The jump loci are closed analytic subsets, by semicontinuity.
\item In rational cohomology $c_1=c_2=0$, and $[\mathrm{ch}(T_X)\,\mathrm{td}(X)]_3=\tfrac12c_1^3-\tfrac{19}{24}c_1c_2+\tfrac12c_3$ (recomputed from Chern roots). Hence $\chi(T_X\otimes M)=\tfrac12c_3=-h^1\le0$.
\end{itemize}
Two points were found in the audit:
\begin{itemize}
\item \emph{Missed generality.} The hypotheses $a(X)=1$ and the holomorphic algebraic reduction enter the proof only to produce a nonzero effective divisor $D$. The note's Lemma~5.1, on which the proof rests, already assumes only such a divisor.
\begin{itemize}
\item \emph{Integral cohomology.} When $H^1(X;\ZZ)=H^2(X;\ZZ)=0$, the bundle $\mathcal O_X(D)$ lies in $\operatorname{Pic}(X)=H^1(X,\mathcal O_X)$ automatically, so the same proof applies to every such threefold that carries a nonzero effective divisor.
\item \emph{Betti numbers.} It suffices that $b_1=b_2=0$. Then $\operatorname{Pic}^0(X)=H^1(X,\mathcal O_X)$ has finite index in $\operatorname{Pic}(X)$, because $H^1(X;\ZZ)=0$ and $H^2(X;\ZZ)$ is finite. Replace $D$ by a multiple $rD$ with $\mathcal O_X(rD)\in\operatorname{Pic}^0(X)$, and work in the component $L_0\cdot\operatorname{Pic}^0(X)$. Rationally $c_1=c_2=0$, and $c_3$ is the Euler number $2-b_3$ \cite[Cor.~11.12 and §14]{MS74}. The conclusion becomes generic vanishing in that component, together with $b_3\ge2$.
\item \emph{A consequence.} Let $X$ be a compact complex threefold with $b_1=b_2=b_3=0$, for example one homeomorphic to $S^6$, that carries a nonzero effective divisor. Then $H^2(X,T_X\otimes L)\neq0$ for every $L\in\operatorname{Pic}(X)$. By the note's Corollary~5.3, also $R^2f_*(T_X\otimes L)\neq0$ for every $L$ and every holomorphic map $f$ onto a curve. So, granting $c_3(X)=2$ and a divisor, both vanishing hypotheses fail for the candidate by this theorem alone. The record's conductor calculation adds where, and by how much, the relative one fails.
\end{itemize}
\item \emph{A citation.} The Riemann--Roch step is applied to $X$, which is not projective (a projective threefold has algebraic dimension $3$). The note cites Hirzebruch's book \cite{Hir66}, which proves the formula for projective manifolds. For compact complex manifolds in general it is a consequence of the Atiyah--Singer index theorem \cite{AS68}.
\end{itemize}
For the constructed threefold, which the record says has $c_3(X)=2$, the contrapositive gives $H^2(X,T_X\otimes L_0)\neq0$ for every $L_0$. This is the use the record makes of the theorem.

\status{\textsf{Audited}. The generalizations are \textsf{proved here}; the first was found by claude-ab, the second and the consequence in the first referee pass.}

\subsection{Relative differentials on the local models}\label{subsec:reldiff}

The record corrects a remark of the source manuscript on the relative differentials of $f:X\to\Pp^1$ (the third \textsf{REFUTED} row, Section~\ref{sec:ledger}). The note's Proposition~6.1 computes three local models. The following statement contains them.

\begin{proposition}[relative differentials of a function on a smooth germ]\label{prop:torsion}
Let $R=\C\{x_1,\dots,x_n\}$ with $n\ge2$, let $0\neq t$ lie in the maximal ideal $\mathfrak m$, and let $M_t=\Omega^1_{R/\C\{t\}}=\bigl(\bigoplus_iR\,dx_i\bigr)/R\,dt$.
\begin{enumerate}[label=\textup{(\alph*)}]
\item Let $\varphi=\gcd(\partial_1t,\dots,\partial_nt)$; $R$ is a unique factorization domain \cite[Thm.~20.3]{Mat86}. Then $\operatorname{Tors}M_t=R\cdot[dt/\varphi]\cong R/(\varphi)$.
\item $M_t$ is free iff $dt(0)\neq0$. Otherwise it needs $n$ generators while its rank is $n-1$.
\item If $t=u\,x_1^{m_1}\cdots x_k^{m_k}$ with $u$ a unit, $1\le k\le n$ and all $m_i\ge1$, then $\varphi=x_1^{m_1-1}\cdots x_k^{m_k-1}$ up to a unit. In particular:
\begin{itemize}
\item for $t=xy$ and $t=xyz$, $M_t$ is torsion-free and not free at the origin;
\item for $t=s^m$ with $m\ge2$, $\operatorname{Tors}M_t=R/(s^{m-1})\,ds$.
\end{itemize}
\item Let $n=3$ and $\Omega^2_{R/\C\{t\}}=\Omega^2_R/(dt\wedge\Omega^1_R)$.
\begin{itemize}
\item For $t=xy$ its torsion is $R/(x,y)\cdot dx\wedge dy\neq0$, supported on the double curve.
\item For $t=xyz$ its torsion is $R\zeta_1+R\zeta_3\cong R^2/\langle(x,0),(y,-y),(0,z)\rangle$, where $\zeta_1=z\,dx\wedge dy+y\,dx\wedge dz$ and $\zeta_3=y\,dx\wedge dz+x\,dy\wedge dz$. It is $\operatorname{Ext}^2_R(R/J,R)$ for the ideal $J=(yz,xz,xy)$ of the three coordinate axes, that is, the canonical module of that curve germ, and it needs two generators.
\item For $t=s^m$ with $m\ge2$ its torsion is $(R/(s^{m-1}))^2$, on $ds\wedge dy_1$ and $ds\wedge dy_2$.
\end{itemize}
\end{enumerate}
\end{proposition}

\begin{proof}
(a) Put $w=\nabla t/\varphi$. Then $\varphi[w]=[dt]=0$. Also $c[w]=0$ iff $cw=d\,\nabla t=d\varphi w$ for some $d$, that is, iff $c\in(\varphi)$; so $R[w]\cong R/(\varphi)$. Conversely, let $f[a]=0$ with $f\neq0$, so $fa=h\,\nabla t=h\varphi w$ for some $h\in R$. Write $h\varphi/f=\alpha/\beta$ in lowest terms. Then $\beta a_i=\alpha w_i$ for all $i$, so $\beta$ divides every $w_i$ and therefore divides $\gcd(w)=1$. So $\beta$ is a unit, and $a\in Rw$. (The same argument shows that $R^n/Rv$ is torsion-free for every $v\in R^n$ whose entries have no common prime factor.)

(b) If some $\partial_it(0)\neq0$, then $dt$ is part of a basis and $M_t\cong R^{n-1}$. If all $\partial_it\in\mathfrak m$, then $M_t/\mathfrak mM_t\cong\C^n$, while $M_t\otimes\operatorname{Frac}R$ has dimension $n-1$.

(c) For $j\le k$, $\partial_jt=x^{m-e_j}\bigl(x_j\partial_ju+m_ju\bigr)$, where the bracket is a unit. For $j>k$, $\partial_jt=x^m\partial_ju$. So every partial is divisible by $\prod_{i\le k}x_i^{m_i-1}$. After dividing by it, the quotient for $j\le k$ is a unit times $\prod_{i\le k,\,i\ne j}x_i$. When $k=1$ this quotient is a unit. When $k\ge2$, a prime dividing all the quotients divides the one for $j=1$, so it is an associate of some $x_i$ with $2\le i\le k$; but the quotient for $j=i$ is not divisible by $x_i$. So the quotients have no common prime factor.

(d) In the basis $dx\wedge dy,\ dx\wedge dz,\ dy\wedge dz$, the relations are $dt\wedge dx_i$.
\begin{itemize}
\item For $t=xy$ they are $(-x,0,0)$, $(y,0,0)$ and $(0,y,x)$. So $\Omega^2_{R/\C\{t\}}=R/(x,y)\oplus R^2/R(y,x)$. The second summand is torsion-free by the argument of (a), since $x$ and $y$ have no common prime factor, and the first is killed by $x$.
\item For $t=xyz$, the map $\omega\mapsto dt\wedge\omega$ sends $\Omega^2_{R/\C\{t\}}$ onto $J\,dx\wedge dy\wedge dz$, which is torsion-free. So the torsion is $\ker(dt\wedge)/\operatorname{im}(dt\wedge)$ on $\Omega^2_R$, once this quotient is known to be killed by $J$, which contains nonzerodivisors. Now $dt\wedge(a\,dx\wedge dy+b\,dx\wedge dz+c\,dy\wedge dz)=(xya-xzb+yzc)\,dx\wedge dy\wedge dz$. If $xya-xzb+yzc=0$, then $x\mid c$; writing $c=xc'$ gives $y(a+zc')=zb$, so $z\mid a$; writing $a=za'$ gives $b=y(a'+c')$. Hence the kernel is free on $\zeta_1=(z,y,0)$ and $\zeta_3=(0,y,x)$. In this basis the relations $dt\wedge dx$, $dt\wedge dy$, $dt\wedge dz$ are $-x\zeta_1$, $y(\zeta_1-\zeta_3)$ and $z\zeta_3$, which gives the presentation. The same computation shows that $0\to R^2\xrightarrow{(\zeta_1\ \zeta_3)}R^3\xrightarrow{(xy,\,-xz,\,yz)}R\to R/J\to0$ is exact. The cokernel of the transpose of the first map, $R^2/\langle(z,0),(y,y),(0,x)\rangle$, is therefore $\operatorname{Ext}^2_R(R/J,R)$, and $(u,v)\mapsto(v,-u)$ identifies it with the presentation above. So the torsion is killed by $J$. Since $R/J$ is Cohen--Macaulay (a reduced curve germ), the module $\operatorname{Ext}^2_R(R/J,R)$ is its canonical module \cite[§3.3]{BH93}. Its relations lie in $\mathfrak m$, so it needs two generators.
\item For $t=s^m$ the relations are $ms^{m-1}ds\wedge dy_1$ and $ms^{m-1}ds\wedge dy_2$.
\end{itemize}
\end{proof}

\paragraph{The global form for one-forms.} This was found in the first referee pass. Let $f:X\to B$ be a holomorphic map from a complex manifold onto a curve, whose fibres are locally $\{u\,x_1^{m_1}\cdots x_k^{m_k}=0\}$. Put $Z=\sum_c\bigl(f^*c-(f^*c)_{\mathrm{red}}\bigr)$. Then
\[
\operatorname{Tors}\Omega^1_{X/B}\;\cong\;\bigl(f^*\omega_B\otimes\mathcal O_X(Z)\bigr)\big|_Z .
\]
\emph{Proof.}
\begin{itemize}
\item Locally $Z$ is cut out by $\varphi_0=\prod_ix_i^{m_i-1}$, which divides $dt$ by Proposition~\ref{prop:torsion}(c).
\item So $dt\otimes h\mapsto h\,dt$ is a holomorphic map $f^*\omega_B\otimes\mathcal O_X(Z)\to\Omega^1_X$.
\item Followed by the projection to $\Omega^1_{X/B}$, its image is the torsion $R\,[dt/\varphi_0]$. Its kernel is $f^*\omega_B$, since $h\,dt\in\mathcal O_X\,dt$ forces $h$ to be holomorphic.
\end{itemize}
For the record's threefold, granting its local models, the multiple fibres $3S_1$ and $4S_2$ give $Z=2S_1+3S_2$.

\paragraph{Consequences for the third \textsf{REFUTED} row.}
\begin{itemize}
\item The record's charts at the cusp fibre are $t=xy$ along its double curves and $t=xyz$ at its triple points. On these charts, Proposition~\ref{prop:torsion}(c) shows that $\Omega^1_{X/B}$ is torsion-free and not locally free, as the record says. The manuscript asserts torsion of $\Omega^1_{X/B}$ along all three singular fibres in its Remark~9.21, so the record's \textsf{REFUTED} label is correct for $p=1$.
\item The manuscript's Remark~1.3, the one named in the ledger, speaks of the sheaves $\Omega^p_{X/B}$. For $p=2$, Proposition~\ref{prop:torsion}(d) shows that the same charts have torsion along the double curves of the cusp fibre, and the multiple-fibre charts have torsion too. So for $p=2$ the torsion statement holds at all three singular fibres.
\item Scope. These are modules on the smooth total space, not restrictions to a fibre (as the note's Remark~6.2 stresses). The multiple fibres enter through the record's description of their fillings as free quotients, with étale covers on which $t=s^m$; that description was not re-checked here.
\end{itemize}

\status{Proposition~\ref{prop:torsion}: \textsf{proved here}. It contains the note's Proposition~6.1, which is \textsf{audited}. Part (d) is new relative to the record; the note's Remark~6.2 leaves exterior powers open.}

\section{The later supplements, 5--9 September}\label{sec:later}

The two later versions of the record add work that uses the construction's data as stated input. Both say that they do not by themselves establish the construction or a counterexample to \cite{CDP20} (guide \fn{14\_}, its opening paragraph and §4; key-advances reader \fn{26\_}, closing paragraph). This section maps them, reproduces their finite statements, and proves one arithmetic statement in full.

\subsection{The topology supplement of 5 September (version 1.2)}

The supplement consists of two papers, \emph{The higher torus wrap} (\fn{15\_}, 50 pages) and \emph{Real $S^{12}$ and $S^{24}$ routes} (\fn{16\_}, 64 pages), and a three-page reader's guide (\fn{14\_}). All three concern maps out of the cusp region of $X$:
\begin{enumerate}[label=\arabic*.]
\item \emph{The angle-forgetting map.} A map $g:X\to S^3$ built from the cusp quotient is based null-homotopic. The proof exhibits an explicit framed null-bordism of the inverse image, a normally framed $T^3$, inside one cusp chart times an interval. It does not use any identification of $X$ with a sphere.
\item \emph{The retained-angle map.} For $H:X\to S^4$ and the reconstruction's fixed diffeomorphism $d:S^6\to X$, the record proves $[H\circ d]=\eta_4\circ\eta_5\neq0$ in $\pi_6(S^4)\cong\ZZ/2$. The inverse image is a framed $T^2$ with the Lie framing. Its spin quadratic form $q(x,y)=x+xy+y$ over $\Ff_2$ has Arf invariant $1$, since $q$ takes the value $1$ on three of the four vectors. This uses the global sphere identification as input.
\item \emph{A Hopf morphism.} Smashing with the Hopf map, $x\mapsto\eta_2\wedge x$, gives an injective homomorphism $\pi_6(S^4)\to\pi_9(S^6)\cong\ZZ/24$ sending $[H\circ d]$ to $12\nu_6$. This is the classical relation $\eta^3=12\nu$ \cite{Toda62} evaluated on the class of item~2, as the guide itself says.
\item \emph{A peripheral morphism.} With the peripheral character $\chi:SL_2(\ZZ)\to\ZZ/12$ and the values $(\chi(C_0),\chi(C_1),\chi(C_2'))=(-1,4,-3)$, the map $\Psi_H=12\chi\bmod24$ takes the values $(12\nu_6,0,12\nu_6)$, and its kernel is $\chi^{-1}(2\ZZ/12)$.
\item \emph{Higher tori.} Over the cusp there is a local map of complex three-tori onto the two-tori, with elliptic kernel, which extends over the toric fillings. The papers do not construct a global higher-torus family. The direct rank-three cusp would have Euler characteristic $6$, and with two Euler-zero finite fibres the total would be $6$, not $\chi(S^8)=2$.
\item \emph{Octonionic cones.} The equal-diagonal positive cone of the octonionic Jordan algebra, compactified at one point, is the pair $(B^{25},S^{24})$. The papers relate it to the flag manifold $F_4/\mathrm{Spin}(8)$; neither completes an analytic torus gluing.
\end{enumerate}
\status{Items 1--3, 5 (the local map) and 6: \textsf{stated}; their proofs are in \fn{15\_} and \fn{16\_}, which were read only in their front matter. \textsf{Reproduced}: the Arf invariant in item~2; the values and kernel in item~4; the Euler-characteristic arithmetic in item~5. The groups $\pi_6(S^4)\cong\ZZ/2$ (generated by $\eta_4\eta_5$) and $\pi_9(S^6)\cong\ZZ/24$, and $\eta^3=12\nu$, are classical \cite{Toda62}.}

\subsection{The key-advances reader of 6 September}

The five-page reader \fn{26\_} gives formulas and proof locations for five results. The proofs are in the frozen project archive (\fn{28\_}; its 782-page workbench and 138-page higher-rung paper), which was not downloaded here.
\begin{enumerate}[label=\textsf{S6-\arabic*}, leftmargin=3.4em]
\item \emph{The anticanonical ring recovers the fibration.} With $f^*(p_1)=3S_1$ and $f^*(p_2)=4S_2$, the reader states $\omega_X\cong\mathcal O_X(-2S_2)$, $\omega_X^{-2}\cong f^*\mathcal O_B(p_2)$ and $f_*\mathcal O_X(kS_2)=\mathcal O_B(\lfloor k/4\rfloor p_2)$. Hence $\bigoplus_mH^0(X,\omega_X^{-m})\cong\C[U,V]$ with $\deg U=1$ and $\deg V=2$, and the degree-two part recovers $f$.
\item \emph{A nonzero deformation.} Replacing $\beta$ by $\beta+c$, with $\operatorname{Im}c$ below an explicit bound, gives a proper holomorphic submersion over a disc of values of $c$. Its Kodaira--Spencer class at each point is nonzero.
\item \emph{The finite fillings.} There are explicit product covers of degrees nine and eight of the central surfaces of the two multiple fibres, and the kernel of the attachment map is an infinite cyclic central subgroup.
\item \emph{Niemeier--Leech arithmetic.} Take the Niemeier lattice $N$ with root system $A_5^4D_4$ and glue index $72$. A neighbour construction with the Weyl vector gives three neighbours, which are cycled by an order-three action. Their common intersection $J$ has determinant $81$ and index $9$ in each, and the fixed extension has root system $A_2^{12}$ with $72$ roots. The value groups of an octonionic cubic on these lattices are computed, with inclusions of index $4$ and $3$.
\item \emph{A common lattice in the triality space.} In Wilson's octonionic Leech construction \cite{Wilson09}, the intersection of the Leech copy with the transported Niemeier lattices is $I=\{\sqrt2(a,a,a):a\in D_4\}$, with Gram matrix $6G_{D_4}$, determinant $5184$ and minimum $12$. On it the cubic is $\tau(\sqrt2(a,a,a))=2\sqrt2\,P(a)$ with $P(a)=a_0\bigl(a_0^2-3(a_1^2+a_2^2+a_3^2)\bigr)$. It generates the group $4\sqrt2\,\ZZ$ but does not take the value $12\sqrt2$.
\end{enumerate}

\begin{proposition}[the arithmetic of S6-5]\label{prop:s65}
Let $D_4=\{a\in\ZZ^4:a_0+a_1+a_2+a_3\equiv0\pmod2\}$, let $a$ run over $D_4$, viewed as quaternions $a_0+a_1i+a_2j+a_3k$, and put $\tau(q)=\operatorname{Re}\bigl((qq)q\bigr)$.
\begin{enumerate}[label=\textup{(\alph*)}]
\item $\tau(\sqrt2\,a)=2\sqrt2\,P(a)$.
\item $P(a)$ is even for every $a\in D_4$, and $P(1,1,0,0)=-2$. So the values $\tau(\sqrt2\,a)$ generate $4\sqrt2\,\ZZ$.
\item For every $a\in\ZZ^4$, if $3\mid P(a)$ then $9\mid P(a)$. So $P$ never takes a value $3k$ with $3\nmid k$. In particular, since $\tau(\sqrt2\,a)=12\sqrt2\,k$ means $P(a)=6k$, $\tau(\sqrt2\,a)\neq12\sqrt2\,k$ whenever $3\nmid k$; the key-advances reader states the case $k=1$.
\item The set of values is exactly $P(D_4)=\{n\in2\ZZ:\ 3\nmid n\text{ or }9\mid n\}$. So $\tau(\sqrt2\,a)=12\sqrt2\,k$ is attained iff $3\mid k$; for example $\tau=36\sqrt2$ at $a=(-3,0,1,2)$.
\end{enumerate}
\end{proposition}

\begin{proof}
(a) Write $a=a_0+v$ with $v$ purely imaginary, so $v^2=-|v|^2$. Then $a^2=a_0^2-|v|^2+2a_0v$, and
\[
\operatorname{Re}(a^3)=a_0(a_0^2-|v|^2)+2a_0\operatorname{Re}(v^2)=a_0(a_0^2-3|v|^2).
\]
Multiply by $(\sqrt2)^3=2\sqrt2$. The quaternions form a subalgebra of the octonions, so the octonionic product agrees.

(b) If $a_0$ is even, so is $P(a)$. If $a_0$ is odd, then $a_1+a_2+a_3$ is odd, so $|v|^2\equiv a_1+a_2+a_3\equiv1$ and $a_0^2-3|v|^2\equiv1-3\equiv0\pmod2$.

(c) $P(a)\equiv a_0^3\pmod3$. So $3\mid P(a)$ forces $a_0=3b$, and then $P(a)=9b(3b^2-|v|^2)$.

(d) By (b) and (c), every value lies in the set. Conversely, write $N=a_1^2+a_2^2+a_3^2$; the condition $a\in D_4$ is $a_0\equiv N\pmod2$. Take an even $n$ in the set.
\begin{itemize}
\item \emph{If $3\nmid n$}, put $k=n/2$, $a_0=-k$ and $N=(k^2+2)/3$. This $N$ is an integer since $3\nmid k$, and then $P=-k(k^2-k^2-2)=n$.
\begin{itemize}
\item If $k$ is odd, then $k^2\equiv1\pmod{24}$, so $N\equiv1\pmod8$.
\item If $k$ is even, then $N\equiv2\pmod4$.
\end{itemize}
\item \emph{If $n=18r$}, put $a_0=-3r$ and $N=3r^2+2$; then $P=-3r(9r^2-9r^2-6)=n$.
\begin{itemize}
\item If $r$ is odd, then $N\equiv5\pmod8$.
\item If $r$ is even, then $N\equiv2\pmod4$.
\end{itemize}
\end{itemize}
In each case $N\equiv a_0\pmod2$, and $N$ is not of the form $4^j(8m+7)$, so it is a sum of three squares by the theorem of Legendre and Gauss \cite[Ch.~IV, Appendix]{Ser73}. For $n=18$ one gets $a_0=-3$ and $N=5=0^2+1^2+2^2$.
\end{proof}

\status{S6-1 to S6-4: \textsf{stated}; the proofs are in the archive \fn{28\_}, which was not downloaded. \textsf{Reproduced}:
\begin{itemize}
\item S6-1: the dimension count $h^0(\omega_X^{-m})=h^0(\Pp^1,\mathcal O(\lfloor m/2\rfloor))=\lfloor m/2\rfloor+1$, which matches the graded pieces of $\C[U,V]$. Also $\omega_X\cong\mathcal O_X(-2S_2)$ agrees with the manuscript's $K_X\cong f^*\mathcal O_B(-1)\otimes\mathcal O_X(2S_2)$, because $f^*\mathcal O_B(1)=\mathcal O_X(4S_2)$.
\item S6-2: $\det_{\R}\Pi=\operatorname{Im}\tau\cdot D$, so the determinant ratio $(D+\operatorname{Im}\delta)/D$ holds.
\item S6-4: the index $72=\sqrt{6^4\cdot4}$; the $72$ roots of $A_2^{12}$; the determinant $81=9^2$.
\item S6-5: the Gram determinant, the minimum, and the formula for $\tau$.
\end{itemize}
Proposition~\ref{prop:s65}: \textsf{proved here}; part (d) was found in the first referee pass. Parts (b) and (c) are also in the paper on four workbenches \cite[Prop.~4.7]{WorkbenchReader}, which took the formula for $\tau$ from the source; part (a) checks that formula. The mod-$3$ step of (c) is the key-advances reader's own argument.}

\section{The integral monodromy of the period system}\label{sec:monodromy}

The integral period system is the part of the construction that can be determined completely by algebra, and it is used elsewhere: the author's higher-torus update attaches a character to its peripheral monodromy (Section~\ref{sec:later}), and the Erdős--Straus workbench builds covers of the projective line from its dual action (Section~\ref{sec:overlaps}). This section determines the monodromy group exactly and explains how these two uses fit it. All results of this section are new, and each is \textsf{proved here}; they were checked by computer and re-derived by an independent referee (Appendix~\ref{app:verification}). The existence question for the complex structure on $S^6$ itself is not addressed.

\subsection{Setting}

The period system is taken from the manuscript \cite{AlpogeS6}, as reconstructed in the record's file \fn{03\_}, §2. Let $V=\ZZ^4$ with basis $(\gamma,u,w,\delta)$. The local monodromies at the two elliptic points are, by formula (2.2) of the manuscript,
\begin{align*}
T_1:&\quad \gamma\mapsto\gamma,\quad u\mapsto-u-w,\quad w\mapsto-6\gamma+u,\quad \delta\mapsto2\gamma+u+w+\delta,\\
T_2:&\quad \gamma\mapsto\gamma,\quad u\mapsto6\gamma+w,\quad w\mapsto-u,\quad \delta\mapsto-3\gamma+u+\delta,
\end{align*}
and the monodromy at the cusp is $T_0=(T_1T_2)^{-1}$. By the manuscript's Lemma~2.2, $T_1^3=T_2^4=I$, and $T_0=I+N$ with $N^2=0$, $N\gamma=Nu=0$, $Nw=-u$ and $N\delta=\gamma$, so $\ker N=\operatorname{im}N=\langle\gamma,u\rangle$.

By the manuscript's Lemma~2.8 the group of $\langle T_1,T_2\rangle$-invariant alternating forms is infinite cyclic, generated by $Q_0$ with $Q_0(\gamma,\delta)=1$ and $Q_0(u,w)=6$, all other pairings of basis vectors being zero. Its elementary divisors are $(1,1,6,6)$, so $Q_0$ has type $(1,6)$. We write $G=\langle T_1,T_2\rangle$ for the monodromy group and
\[
G_{\ZZ}=\{g\in Sp(Q_0,\ZZ):\ g\gamma=\gamma\}
\]
for the stabiliser of $\gamma$. Every $T_i$ fixes $\gamma$ and preserves $Q_0$, so $G\subseteq G_{\ZZ}$.

\subsection{The stabiliser of \texorpdfstring{$\gamma$}{gamma}}

For $M\in SL_2(\ZZ)$ let $L_M$ act by $M$ on the plane $\langle u,w\rangle$, in the basis $(u,w)$, and fix $\gamma$ and $\delta$. For $p,r,s\in\ZZ$ let $h(p,r,s)$ be the map
\[
\gamma\mapsto\gamma,\qquad u\mapsto u-6r\gamma,\qquad w\mapsto w+6p\gamma,\qquad \delta\mapsto\delta+pu+rw+s\gamma.
\]

\begin{lemma}\label{lem:stab}
Every element of $G_{\ZZ}$ can be written uniquely as $L_Mh(p,r,s)$. The elements $h(p,r,s)$ form a group $\mathrm{Heis}(\ZZ)$ with the law
\[
h(p_1,r_1,s_1)\,h(p_2,r_2,s_2)=h\bigl(p_1+p_2,\ r_1+r_2,\ s_1+s_2+6(p_1r_2-p_2r_1)\bigr),
\]
and $L_Mh(p,r,s)L_M^{-1}=h(M(p,r),s)$. Hence $G_{\ZZ}=SL_2(\ZZ)\ltimes\mathrm{Heis}(\ZZ)$.
\end{lemma}
\begin{proof}
Let $g\in G_{\ZZ}$. Since $g$ fixes $\gamma$, it preserves $\gamma^\perp=\langle\gamma,u,w\rangle$. The form induced on $\gamma^\perp/\langle\gamma\rangle$ is $6\,(u^*\wedge w^*)$, so the induced map $M$ has determinant $1$. On $V/\gamma^\perp$, $g$ acts trivially, since $Q_0(\gamma,gx)=Q_0(g^{-1}\gamma,x)=Q_0(\gamma,x)$. Then $L_M^{-1}g$ acts trivially on $\langle\gamma\rangle$, on $\gamma^\perp/\langle\gamma\rangle$ and on $V/\gamma^\perp$. Such a unipotent map has the form $u\mapsto u+a_u\gamma$, $w\mapsto w+a_w\gamma$, $\delta\mapsto\delta+pu+rw+s\gamma$, and it preserves $Q_0$ exactly when $a_u=-6r$ and $a_w=6p$. The law and the conjugation formula are matrix products.
\end{proof}

Two consequences are used repeatedly. Inside $\mathrm{Heis}(\ZZ)$, conjugation acts by
\[
h(p_0,r_0,s_0)\,h(p,r,s)\,h(p_0,r_0,s_0)^{-1}=h\bigl(p,r,s+12(p_0r-pr_0)\bigr).
\] The centre of $G_{\ZZ}$ is $\{h(0,0,s)\}=\langle t_\gamma\rangle$, where $t_\gamma=h(0,0,1)$ is the transvection $\delta\mapsto\delta+\gamma$, that is, $x\mapsto x+Q_0(\gamma,x)\gamma$. The subgroup $H(\ZZ)=\{h(p,r,6\kappa)\}$ carries the standard Heisenberg law $\kappa_1+\kappa_2+(p_1r_2-p_2r_1)$ and has index $6$ in $\mathrm{Heis}(\ZZ)$; so $G_{\ZZ}$ contains the Jacobi group $SL_2(\ZZ)\ltimes H(\ZZ)$ with index $6$, and is generated by it together with $t_\gamma$.

\subsection{The monodromy group is the kernel of a character of order 12}

Let $\operatorname{ab}:SL_2(\ZZ)\to\ZZ/12$ be the abelianisation, normalised by $\operatorname{ab}(T)=1$ for $T=\begin{psmallmatrix}1&1\\0&1\end{psmallmatrix}$. Then $\operatorname{ab}(S)=9$ for $S=\begin{psmallmatrix}0&-1\\1&0\end{psmallmatrix}$, $\operatorname{ab}(R)=4$ for $R=\begin{psmallmatrix}-1&1\\-1&0\end{psmallmatrix}$, and $\operatorname{ab}(-I)=6$. These values are forced: $SL_2(\ZZ)^{\mathrm{ab}}\cong\ZZ/12$ is generated by the class of $T$ (the record's file \fn{15\_} proves this by a Smith normal form in its Theorem~6.1); $S$ and $R$ have orders $4$ and $3$; and $\operatorname{ab}(S)+\operatorname{ab}(R)=\operatorname{ab}(SR)=\operatorname{ab}(T)$, because $SR=\begin{psmallmatrix}1&0\\-1&1\end{psmallmatrix}$ is conjugate to $T$. On $\mathrm{Heis}(\ZZ)$ put
\[
\varphi(h(p,r,s))=s-6q(p,r)\bmod12,\qquad q(p,r)=p^2+pr+r^2\bmod2,
\]
and on $G_{\ZZ}$ put $\Phi(L_Mh)=\operatorname{ab}(M)+\varphi(h)\bmod12$.

\begin{theorem}[the monodromy group]\label{thm:mono}
\begin{enumerate}[label=(\alph*)]
\item $\Phi:G_{\ZZ}\to\ZZ/12$ is a surjective homomorphism, and $G=\ker\Phi$.
\item The local monodromies decompose as $T_1=L_Rh(-1,0,2)$, $T_2=L_Sh(0,-1,-3)$ and $T_0=L_{T^{-1}}t_\gamma$. In particular the Levi image of $G$ is all of $SL_2(\ZZ)$.
\item $G$ is normal in $G_{\ZZ}$, and $G_{\ZZ}/G\cong\ZZ/12$ is cyclic, generated by the class of the central transvection $t_\gamma$. Moreover $G_{\ZZ}=G\cdot\langle t_\gamma\rangle$ and $G\cap\langle t_\gamma\rangle=\langle t_\gamma^{12}\rangle$.
\item $G\cap\mathrm{Heis}(\ZZ)=H':=\{h(p,r,s):s\equiv6(p^2+pr+r^2)\bmod12\}$. This subgroup has index $12$, with coset representatives $h(0,0,k)$, $0\le k\le11$, and it is generated by $c=[T_1,T_2^2]=h(1,0,-6)$ and $c'=T_2cT_2^{-1}=h(0,1,6)$.
\item $G\cap SL_2(\ZZ)=SL_2(\ZZ)'$, the commutator subgroup, which is free of rank $2$ and of index $12$.
\item The abelianisation of $G_{\ZZ}$ is $(\ZZ/12)^2$, via $g\mapsto(\operatorname{ab}(M),\varphi(h))$, and $[G_{\ZZ},G_{\ZZ}]=SL_2(\ZZ)'\ltimes H'$.
\item $G$ contains the principal congruence subgroup of level $12$ and no principal congruence subgroup of level $4$ or $6$. For every $D\ge1$ the image of $G$ in $GL_4(\ZZ/D)$ has index $\gcd(D,12)$ in the image of $G_{\ZZ}$.
\item $G$ is Zariski dense in $P=SL_2\ltimes\mathrm{Heis}$, the algebraic stabiliser of $\gamma$, and it is an arithmetic subgroup of $P$, since it has finite index in $P(\ZZ)=G_{\ZZ}$.
\end{enumerate}
\end{theorem}
\status{Proved here; new. Checked by computer (\note{s6\_monodromy\_checks.py}, $24$ checks; \note{ver52\_referee\_claims.py}, $22$ checks, including that every one of the $165{,}888$ elements of $G$ modulo $12$ has $\Phi=0$), and confirmed by the referee's constructive membership algorithm.}
\begin{proof}
\emph{(a), $\Phi$ is a homomorphism.} $\varphi$ is a homomorphism on $\mathrm{Heis}(\ZZ)$ because $q$ is a quadratic refinement of the mod-$2$ symplectic form $\omega(x,y)=p_xr_y-p_yr_x$: $q(x+y)\equiv q(x)+q(y)+\omega(x,y)\pmod2$, and therefore $6q(x+y)\equiv6q(x)+6q(y)+6\omega(x,y)\pmod{12}$. Moreover $\varphi$ is invariant under the Levi action, since $q$ equals $1$ exactly on the three nonzero vectors of $\Ff_2^2$ and so $q(Mx)\equiv q(x)$. Writing $L_{M_1}h_1\cdot L_{M_2}h_2=L_{M_1M_2}\cdot(L_{M_2}^{-1}h_1L_{M_2})\cdot h_2$ gives $\Phi(g_1g_2)=\Phi(g_1)+\Phi(g_2)$.

\emph{(a), $G\subseteq\ker\Phi$ and a lower bound.} By (b), $\Phi(T_1)=4+(2-6)=0$ and $\Phi(T_2)=9+(-3-6)=0$, so $G\subseteq\ker\Phi$. Since $\Phi(t_\gamma)=1$, $\ker\Phi$ has index $12$; hence $[G_{\ZZ}:G]\ge12$ and $G\cap\mathrm{Heis}(\ZZ)\subseteq\ker\varphi=H'$.

\emph{(a), $G=\ker\Phi$.} The words $c^pc'^r[c,c']^m$ equal $h(p,r,-6p+6r+6pr+12m)$, and $-6p+6r+6pr\equiv6q(p,r)\pmod{12}$, so these words exhaust $H'$ and $H'=\langle c,c'\rangle\subseteq G$. Since $R$ and $S$ generate $SL_2(\ZZ)$ ($S(SR)S^{-1}=T$), $G$ maps onto $SL_2(\ZZ)$. Therefore $[G_{\ZZ}:G]=[\mathrm{Heis}(\ZZ):G\cap\mathrm{Heis}(\ZZ)]\le[\mathrm{Heis}(\ZZ):H']=12$, and $G=\ker\Phi$.

(b) is a direct computation, and (c) follows from (a) and the description of the centre.

(d) $H'$ is $\ker\varphi$, and $h(0,0,k)h(p,r,s)=h(p,r,s+k)$ gives the coset representatives.

(e) $L_M\in G$ exactly when $\operatorname{ab}(M)=0$. The commutator subgroup $SL_2(\ZZ)'$ has index $12$ and is torsion-free: every element of finite order in $SL_2(\ZZ)$ is conjugate to a power of $S$ or of $-R$, and these powers have the $\operatorname{ab}$-values $9,6,3$ and $10,8,6,4,2$, all nonzero. A torsion-free subgroup of the amalgam $\ZZ/4*_{\ZZ/2}\ZZ/6$ acts freely on its Bass--Serre tree and is free \cite[Ch.~I]{SerreTrees}; its rank follows from the Euler characteristic, $12\cdot(-\frac1{12})=-1=1-\mathrm{rank}$.

(f) The map is a homomorphism by the computation in (a), applied to each coordinate, and it is onto, since $L_T\mapsto(1,0)$ and $t_\gamma\mapsto(0,1)$. Its kernel $SL_2(\ZZ)'\ltimes H'$ is generated by commutators: those of Levi elements give $SL_2(\ZZ)'$; the commutators $[L_M,h(x,s)]=h((M-I)x,-6\omega(Mx,x))$ have projections $(M-I)x$ that span $\ZZ^2$, for example $(T-I)(0,1)=(1,0)$ and $(T^{\mathsf T}-I)(1,0)=(0,1)$; and the commutators of these give $h(0,0,12)$. By Proposition~\ref{prop:heis} below the subgroup they generate is $H'$.

(g) If $g\equiv I\pmod{12}$, then $\operatorname{ab}(M)=0$, because $\operatorname{ab}$ factors through $SL_2(\ZZ/12)$ (checked by a breadth-first search over all $1152$ elements), and $\varphi(h)=0$. On the other hand $h(0,0,4)$ and $h(0,0,6)$ lie in the principal congruence subgroups of levels $4$ and $6$ and have $\Phi=4$ and $6$. For the index, let $K_D$ be the principal congruence subgroup of level $D$; then $[G_{\ZZ}\bmod D:G\bmod D]=[G_{\ZZ}:GK_D]=12/|\Phi(K_D)|$. Since $t_\gamma^D\in K_D$, $\Phi(K_D)$ contains $\gcd(D,12)\ZZ/12$. Conversely, $\Phi$ modulo $d=\gcd(D,12)$ factors through reduction modulo $d$: $SL_2(\ZZ/12)=SL_2(\ZZ/4)\times SL_2(\ZZ/3)$ with $SL_2(\ZZ/3)^{\mathrm{ab}}=\ZZ/3$, and $6q(p,r)\bmod d$ depends only on $p,r\bmod2$.

(h) $h(2,0,0)$, $h(0,2,0)$ and $h(0,0,12)$ lie in $H'\subseteq G$. Their cyclic groups are Zariski dense in the three one-parameter subgroups that generate $\mathrm{Heis}$, so the Zariski closure of $G$ contains $\mathrm{Heis}$. Its Levi image is closed and contains $SL_2(\ZZ)$, which is dense in $SL_2$.
\end{proof}

In particular $G$ is not a thin group in the sense of Sarnak \cite{SarnakThin}: it has finite index in the integral points of its Zariski closure. This places the $S^6$ period system among the arithmetic cases of rank-four symplectic monodromy; Brav and Thomas exhibit hypergeometric monodromy groups in $Sp(4,\ZZ)$ of the other, thin kind \cite{BravThomas}.

\subsection{The Heisenberg part, and the twist}

\begin{proposition}\label{prop:heis}
\begin{enumerate}[label=(\alph*)]
\item The quadratic refinements of $\omega$ modulo $2$ on $\Ff_2^2$ are the four functions $\alpha p^2+pr+\beta r^2$ ($\alpha,\beta\in\Ff_2$). Three have Arf invariant $0$, and $q=p^2+pr+r^2$ has Arf invariant $1$. $SL_2(\Ff_2)$ permutes the three even refinements and fixes $q$.
\item Let $N\subseteq\mathrm{Heis}(\ZZ)$ be a subgroup that is normalised by the Levi group and maps onto $\ZZ^2$ under $(p,r)$. Then $N=N_k:=\{h(p,r,s):s\equiv6q(p,r)\bmod k\}$ for a unique divisor $k$ of $12$. Conversely, every $N_k$ is such a subgroup, of index $k$.
\item $H'=N_{12}$ is the smallest of them, and $H'=[G_{\ZZ},\mathrm{Heis}(\ZZ)]=[G_{\ZZ},G_{\ZZ}]\cap\mathrm{Heis}(\ZZ)$.
\item On $H(\ZZ)=\{h(p,r,6\kappa)\}$, $H'$ is the kernel of the character $v_H(h(p,r,6\kappa))=(-1)^{p+r+pr+\kappa}$.
\end{enumerate}
\end{proposition}
\status{Proved here; new.}
\begin{proof}
(a) is a finite check. (b) Lifts of $(1,0)$ and $(0,1)$ in $N$ have commutator $h(0,0,12)$, so $N$ meets the centre in $k\ZZ$ for some $k\mid12$, and the fibre of $N$ over $x\in\ZZ^2$ is $\sigma(x)+k\ZZ$. The function $\lambda(x)=\sigma(x)-6q(x)\bmod k$ is additive, by the group law and (a), and it is Levi-invariant; so it vanishes on every $(M-I)\ZZ^2$, and these span $\ZZ^2$. Hence $\lambda=0$. (c) The subgroup $[G_{\ZZ},\mathrm{Heis}(\ZZ)]$ is Levi-normal, maps onto $\ZZ^2$, contains $h(0,0,12)$ and lies in $H'$, so it is $N_{12}=H'$. (d) With $s=6\kappa$, the condition $s\equiv6q(p,r)\pmod{12}$ says $\kappa\equiv q(p,r)\equiv p+pr+r\pmod2$.
\end{proof}

Two readings of Proposition~\ref{prop:heis} connect the Heisenberg part to established theory. Quadratic refinements of the mod-$2$ intersection form of a surface are its spin structures, with the Arf invariant as their parity \cite{Johnson}; on a Riemann surface they are its theta characteristics \cite{AtiyahSpin}. On a torus the odd spin structure is the one fixed by the mapping class group $SL_2(\ZZ)$, and (b) shows that $SL_2$-invariance forces it. The character $v_H$ of (d) is the Heisenberg multiplier of the odd Jacobi theta function $\vartheta$: Gritsenko and Hulek show that $\vartheta\in J_{1/2,1/2}(v_\eta^3\times v_H)$ \cite[(1.4)--(1.5)]{GritsenkoHulek}. So the Heisenberg part of the monodromy group is the kernel of the odd theta multiplier.

\begin{proposition}\label{prop:twist}
For $k\in\ZZ/12$ put $G_k=\{L_Mh:\varphi(h)+k\operatorname{ab}(M)\equiv0\}$. The $G_k$ are twelve distinct normal subgroups of index $12$. Each has Levi image $SL_2(\ZZ)$ and Heisenberg part $H'$, and $G=G_1$. The extension $1\to H'\to G\to SL_2(\ZZ)\to1$ does not split, whereas $G_0=SL_2(\ZZ)\ltimes H'$ does.
\end{proposition}
\status{Proved here; new.}
\begin{proof}
Distinctness: $L_Th(0,0,s)\in G_k$ exactly when $s\equiv-k$. Non-splitting: a splitting $\sigma$ would give an involution $\sigma(-I)\in G$ over $-I$ that commutes with an element $\sigma(R)\in G$ over $R$. Since $(L_{-I}h(x,s))^2=h(0,2s)$, the involutions of $G$ over $-I$ are the $L_{-I}h(x,0)$ with $\Phi=6-6q(x)\equiv0$, that is, with $x\not\equiv0\pmod2$. Conjugation by $L_Rh(y,t)$ sends $L_{-I}h(x,0)$ to $L_{-I}h(R(x-2y),0)$, so commuting would require $Rx\equiv x\pmod2$; but $R$ has no nonzero fixed vector on $\Ff_2^2$.
\end{proof}

So the pair \emph{index $12$ in the centre, odd refinement modulo $2$} determines the Heisenberg part $H'$, and the group $G$ itself is the twist of $H'$ by the abelianisation of the Levi part.

\subsection{The marked peripheral character of the higher-torus update}

The higher-torus update of the record (its file \fn{15\_}, §3.2 and Theorem~6.1; see also \fn{14\_}, §2) works with the peripheral matrices
\[
C_1=\begin{pmatrix}-1&-1\\1&0\end{pmatrix},\qquad C_2'=\begin{pmatrix}0&1\\-1&0\end{pmatrix},\qquad C_0=\begin{pmatrix}1&1\\0&1\end{pmatrix}.
\]
It defines the marked peripheral character $\chi=\nu\circ\operatorname{ab}:SL_2(\ZZ)\twoheadrightarrow\ZZ/12$ by $\chi(C_0)=-1$, and shows that $(\chi(C_0),\chi(C_1),\chi(C_2'))=(-1,4,-3)$.

\begin{theorem}[the peripheral character]\label{thm:periph}
$C_1=\pi(T_1)^{\mathsf T}$, $C_2'=\pi(T_2)^{\mathsf T}$ and $C_0=T$, where $\pi$ denotes the Levi part, and $\chi=-\operatorname{ab}$. The triple $(-1,4,-3)$ equals $(\operatorname{ab}\pi(T_0),\operatorname{ab}\pi(T_1),\operatorname{ab}\pi(T_2))$. For every $g=L_Mh(p,r,s)$ in the monodromy group $G$,
\[
\chi(M)\equiv s-6(p^2+pr+r^2)\pmod{12}.
\]
\end{theorem}
\status{Proved here; new.}
\begin{proof}
The matrix identities are direct. $\chi$ and $-\operatorname{ab}$ are homomorphisms that agree on $T$, which generates the abelianisation, so they are equal. The triple follows from Theorem~\ref{thm:mono}(b), and the last formula is $\Phi(g)=0$ rewritten.
\end{proof}

So the character that the update attaches to the peripheral monodromy is exactly the Levi half of the character whose kernel is the monodromy group. On $G$ it is determined by the Heisenberg coordinates, and the monodromy group is the locus where the two halves cancel. The same update exhibits a Lie-framed torus whose spin quadratic form is $q(x,y)=x+xy+y$ over $\Ff_2$, with Arf invariant one (\fn{15\_}, (147); \fn{14\_}, §1); since $x+xy+y\equiv x^2+xy+y^2\pmod2$, this is the odd refinement of Proposition~\ref{prop:heis}(a). Whether that torus's first homology is the lattice $\langle\bar u,\bar w\rangle$ of Proposition~\ref{prop:heis} is Question~\ref{q:arfone}.

\subsection{The orbits of the \texorpdfstring{$(3,4,\infty)$}{(3,4,infinity)} covers at every level}

The Erdős--Straus workbench (its package ES-09) uses the dual matrices $A_j=(T_j^{-1})^{\mathsf T}$, acting on $\Lambda=\operatorname{Hom}(V,\ZZ)$ with dual basis $(\hat\gamma,\hat u,\hat w,\hat\delta)$, to build covers from the orbits on the affine hyperplane
\[
H_D=\{\lambda\in\Lambda/D\Lambda:\ \text{the $\hat\gamma$-coordinate of $\lambda$ is }1\}.
\]
Write $\lambda=\hat\gamma+b\hat u+c\hat w+d\hat\delta$ and $e=\gcd(D,6)$.

\begin{theorem}[the orbits at every level]\label{thm:orbitsS6}
For every $D\ge1$ the orbits of $G$ on $H_D$ coincide with the orbits of $G_{\ZZ}$. The orbit of $\lambda$ is $\{\lambda'\in H_D:\gcd(b',c',e)=\gcd(b,c,e)\}$, of size $D^3\cdot\#\{v\in(\ZZ/e)^2:\gcd(v,e)=\gcd(b,c,e)\}/e^2$.
\end{theorem}
\status{Proved here; new. Checked as full orbit partitions for all $D\le14$ (\note{ver52\_referee\_claims.py}) and by the referee for all $D\le24$.}
\begin{proof}
\emph{The dual action.} The dual action of $h(p,r,s)$ is $(b,c,d)\mapsto(b+6r,\,c-6p,\,d-s-pb-rc)$, and $L_M$ acts on $(b,c)$ by $M^{-\mathsf T}$ and fixes $d$.

\emph{The orbits of $G_{\ZZ}$.} The coordinate $d$ can be moved freely, and $(b,c)$ matters only modulo $6(\ZZ/D)^2$. So the $G_{\ZZ}$-orbits correspond to the $SL_2(\ZZ/e)$-orbits on $(\ZZ/e)^2$, which are classified by content, since $SL_2(\ZZ)$ maps onto $SL_2(\ZZ/e)$.

\emph{$G$ has the same orbits.} $G$ is normal and $G_{\ZZ}=G\langle t_\gamma\rangle$ with $t_\gamma$ central, so it suffices to show $t_\gamma\lambda\in G\lambda$. At a point with $(b,c)=(g,0)$ the Levi element $L_{(T^{\mathsf T})^{-1}}$ fixes $\lambda$ and has $\Phi=1$, because $(T^{\mathsf T})^{-1}=STS^{-1}$; hence $t_\gamma L_{(T^{\mathsf T})^{-1}}^{-1}\in G$ moves $\lambda$ exactly as $t_\gamma$ does. Every point is $G_{\ZZ}$-equivalent to one with $c=0$, and normality and centrality transport the conclusion to every point.
\end{proof}

\begin{corollary}\label{cor:orbitsS6}
For $D=2^k$, $k\ge1$, there are exactly two orbits, of sizes $D^3/4$ and $3D^3/4$. For odd $D$ there is one orbit if $3\nmid D$, and there are two, of sizes $D^3/9$ and $8D^3/9$, if $3\mid D$. For $6\mid D$ there are four orbits; at $D=6$ their sizes are $6$, $18$, $48$ and $144$, and at $D=12$ they are $48$, $144$, $384$ and $1152$. In general the orbits at $D=2^km$, $m$ odd, are the products of those at $2^k$ and at $m$.
\end{corollary}
\status{Proved here; new.}
\begin{proof}
Count the vectors of $(\ZZ/e)^2$ by content, as in the Erdős--Straus paper \cite[Corollary~5.2]{ESPaper}, and multiply by $D^3/e^2$; the product statement is the Chinese remainder theorem.
\end{proof}

This settles the question, raised when these covers were examined at even levels $D=2^k$ with $k\le5$, whether the two-orbit structure persists for all $k$ and whether mixed levels decompose as products.

\subsection{What the period system does not have}\label{ssec:mononeg}

Three statements of this section and of the manuscript are negative in form, and each is proved.
\begin{itemize}
\item \emph{The extension does not split} (Proposition~\ref{prop:twist}). No subgroup of $G$ maps isomorphically onto $SL_2(\ZZ)$: an involution over $-I$ would have to commute with an element over $R$, and $R$ has no nonzero fixed vector on $\Ff_2^2$.
\item \emph{The level is exactly $12$} (Theorem~\ref{thm:mono}(g)). $G$ contains no principal congruence subgroup of level $4$ or $6$, since $h(0,0,4)$ and $h(0,0,6)$ lie in them and have $\Phi=4$ and $6$.
\item \emph{No polarised limit mixed Hodge structure at the cusp} (the manuscript's Remark~3.23). With $N$ as above, the form $b_N(x,y)=Q_0(x,Ny)$ on $V/\ker N=\langle w,\delta\rangle$ has $b_N(w,w)=Q_0(w,-u)=6$, $b_N(\delta,\delta)=Q_0(\delta,\gamma)=-1$ and $b_N(w,\delta)=b_N(\delta,w)=0$. So $b_N=\operatorname{diag}(6,-1)$ is indefinite, and $(V,Q_0,N)$ does not support a polarised limit mixed Hodge structure.
\end{itemize}
These results say which standard tools apply as they stand: methods that need a split extension or a polarised limit do not, while methods for congruence subgroups of level $12$ and for arithmetic groups do.

\section{Bridges to other areas and to the author's other workbenches}\label{sec:overlaps}

In the author's view the connections of this research to other areas are among its most valuable outcomes. This section first treats the connections of the period system to established theories, and then the typed edges to the author's other workbenches. For each it states what is established and what is open.

\subsection{The period system and established theories}\label{ssec:theories}

\paragraph{The paramodular group of level $6$.} Gritsenko and Hulek define, for the alternating form $W_t$ of type $(1,t)$ on $\ZZ^4$, the integral paramodular group $\widetilde\Gamma_t=Sp(W_t,\ZZ)$, and they realise it as a group $\Gamma_t$ of rational symplectic matrices \cite[§1]{GritsenkoHulek}. Two alternating forms on $\ZZ^4$ with the same elementary divisors are equivalent under $GL_4(\ZZ)$ (the Frobenius normal form of alternating integral matrices), and $Q_0$ has elementary divisors $(1,1,6,6)$ (Section~\ref{sec:monodromy}). So $Sp(Q_0,\ZZ)\cong\widetilde\Gamma_6$: the ambient group of the period system is the paramodular group of level $6$. It contains $G_{\ZZ}$ as the stabiliser of the primitive isotropic vector $\gamma$, whose divisor $Q_0(\gamma,V)$ is $\ZZ$, and $G_{\ZZ}$ contains the Jacobi group $SL_2(\ZZ)\ltimes H(\ZZ)$ with index $6$ (Section~\ref{sec:monodromy}).

The character $e(\Phi/12)$, with $e(x)=e^{2\pi ix}$, has the following restrictions:
\begin{itemize}
\item on the Levi group, $e(\operatorname{ab}(M)/12)=v_\eta^2(M)$, where $v_\eta$ is the multiplier of Dedekind's $\eta$. Indeed $v_\eta^2$ is the multiplier of the weight-one form $\eta^2$, hence a character of $SL_2(\ZZ)$, and $v_\eta^2(T)=e(1/12)$ because $\eta(\tau+1)=e(1/24)\eta(\tau)$; two characters of $SL_2(\ZZ)$ that agree on $T$ are equal, since $T$ generates the abelianisation;
\item on the standard Heisenberg group $H(\ZZ)=\{h(p,r,6\kappa)\}$, identified with the integral Heisenberg group by $h(p,r,6\kappa)\mapsto[p,r;\kappa]$, it is $v_H[p,r;\kappa]=(-1)^{p+r+pr+\kappa}$ (Proposition~\ref{prop:heis}(d)), which is the Heisenberg multiplier of the odd Jacobi theta function in \cite[(1.4)--(1.5)]{GritsenkoHulek};
\item on the centre, $e(\Phi(h(0,0,s))/12)=e(s/12)$.
\end{itemize}
For $t=6$, Lemma~1.1 of \cite{GritsenkoHulek} gives characters $\chi_3$ and $\chi_4$, of orders $3$ and $4$, of the extended group $\Gamma_6^+$. On the Jacobi parabolic of their realization they restrict to $v_\eta^{8}$ and $v_\eta^{6}$ on $SL_2(\ZZ)$, to $v_H^{4}=1$ and $v_H^{3}=v_H$ on $H(\ZZ)$, and to $e(\kappa/3)$ and $e(\kappa/4)$ on the central elements $[0,0;\kappa/6]$. Hence $\chi_{12}=\chi_3\chi_4^{-1}$ has order $12$ and restricts there to $v_\eta^2$, $v_H$ and $e(\kappa/12)$: the same three restrictions as $e(\Phi/12)$. The Lemma was read through an extraction of its statement, not in the full text. Whether the parabolic of \cite{GritsenkoHulek} is conjugate in $\Gamma_6^+$ to $G_{\ZZ}$, and whether $e(\Phi/12)$ is the restriction of $\chi_{12}$ to $G_{\ZZ}$, is not verified here (Question~\ref{q:paramodular}). My assessment: I think the identification is likely, because the two characters have the same order and the same restrictions to the three pieces; what has to be computed is the value of $\chi_{12}$ on the two minimal positive transvections $L_T$ (along $u$, of divisor $6$) and $t_\gamma$ (along $\gamma$, of divisor $1$), which both have $\Phi=1$. Since $G_{\ZZ}^{\mathrm{ab}}$ is generated by the classes of these two elements (Theorem~\ref{thm:mono}(f)), if the two values are equal and of order $12$, then $G=G_{\ZZ}\cap\ker\chi_{12}$, and the monodromy group of the $S^6$ period system is cut out of the stabiliser of $\gamma$ in the paramodular group by a character of Gritsenko and Hulek.

\paragraph{Spin structures and theta characteristics.} Quadratic refinements of the mod-$2$ intersection form of a surface are its spin structures, with the Arf invariant as their parity \cite{Johnson}; on a Riemann surface they are its theta characteristics \cite{AtiyahSpin}. On a torus the odd spin structure is the one fixed by the mapping class group $SL_2(\ZZ)$, and Proposition~\ref{prop:heis}(b) shows that $SL_2$-invariance forces it in the monodromy group. The record's higher-torus update exhibits a Lie-framed torus whose spin form $x+xy+y$ has Arf invariant one, and this is the odd refinement (Section~\ref{sec:monodromy}). Whether that torus's first homology is the lattice $\langle\bar u,\bar w\rangle$ is Question~\ref{q:arfone}. My assessment: if it is, the update realises geometrically the spin structure that cuts out the Heisenberg part of the monodromy, and the update's construction and the monodromy computation describe the same object.

\paragraph{The cusp.} The cusp monodromy is $T_0=I+N$ with $N^2=0$ and $\operatorname{rank}N=2$. On $\Lambda^2V$ the induced nilpotent $N_2(a\wedge b)=Na\wedge b+a\wedge Nb$ satisfies $N_2^2(a\wedge b)=2\,Na\wedge Nb$ and $N_2^3=0$, and $N_2^2(w\wedge\delta)=-2\,u\wedge\gamma\neq0$. So $N_2^2\neq0=N_2^3$. For $H^2$ of a semistable degeneration of surfaces with trivial canonical class, $N^2\neq0$ is the monodromy of Kulikov's type~III \cite{Kulikov,PerssonPinkham}. The manuscript describes its central fibre at the cusp as a degree-six del Pezzo surface glued to itself along opposite sides of its hexagon, with two triple points (Section~\ref{sec:candidate}); its dual complex then has one vertex, three edges and two triangles, a triangulation of the torus (Euler characteristic $1-3+2=0$). This is the combinatorial type of Mumford's degenerations of abelian surfaces. The dictionary of Kulikov and of Persson and Pinkham rests on polarised limits, and $(V,Q_0,N)$ supports none, since $b_N$ is indefinite (Section~\ref{ssec:mononeg}). So the comparison holds for the Jordan type and the combinatorics. My assessment: I think the combinatorial match indicates that the degeneration is of Mumford type, and a substitute for the polarisation is what a Hodge-theoretic comparison needs.

\paragraph{The modular part.} Up to the factor $6$ in the induced form, the Levi group acts on $\gamma^\perp/\langle\gamma\rangle$ as $SL_2(\ZZ)$ acts on the first homology of an elliptic curve, and the Heisenberg part records the extension data. Theorem~\ref{thm:mono} and Proposition~\ref{prop:heis} describe both exactly. Whether $\gamma^\perp/\langle\gamma\rangle$ is the homology of an elliptic family inside or under the tori is Question~\ref{q:elliptic}.

\paragraph{Arithmetic, not thin.} $G$ has finite index in the integral points of its Zariski closure (Theorem~\ref{thm:mono}(h)), so it is not thin in the sense of Sarnak \cite{SarnakThin}; Brav and Thomas exhibit hypergeometric monodromy groups in $Sp(4,\ZZ)$ of the thin kind \cite{BravThomas}. Arithmeticity makes the methods for congruence subgroups available, and they give the orbit classification of Theorem~\ref{thm:orbitsS6}.

\subsection{The edges to the author's other workbenches}

The audit note \note{47\_} records every crossing between the owner's workbenches that the audit found, as typed edges: \emph{P}, a proof-level import; \emph{I}, the same identity or object on both sides; \emph{A}, an analogy or a shared number only. The table lists the edges that touch the $S^6$ material.

\begin{center}\small
\begin{longtable}{@{}>{\raggedright\arraybackslash}p{0.13\linewidth}>{\raggedright\arraybackslash}p{0.44\linewidth}>{\raggedright\arraybackslash}p{0.07\linewidth}>{\raggedright\arraybackslash}p{0.28\linewidth}@{}}
\toprule
Edge & What crosses & Type & Status\\
\midrule
\endhead
$S^6\to$ YM & The imaginary part of the period matrix, used as a magnetic background in the Yang--Mills workbench. Only $D_{\mathrm{YM}}=L_Tq_T+6m_T^2=-\det_{\R}\Pi$ enters (Section~\ref{sec:candidate}) & P & audited (\cite{YMPaper}, §§6.3 and~8)\\
$S^6\to$ ES & The integral monodromies $A_j=(T_j^{-1})^{\mathsf T}$ of the $(3,4,\infty)$ period system, in the Erdős--Straus package ES-09, with $A_1^3=A_2^4=A_1A_2A_\infty=I$. Nothing else from the construction is imported & P & the orbits at every level: Theorem~\ref{thm:orbitsS6}; the package's statements are correct for odd levels, which the application uses, and do not hold as stated at even levels \cite[§5.2]{ESPaper}\\
ES $\leftrightarrow$ $S^6$ & The Niemeier lattice with root system $A_5^4D_4$ and glue index $72$: the Erdős--Straus result R10, and S6-4 & I & the Erdős--Straus glue code verified; the $S^6$ code is in the record's archive \fn{28\_}, and the two were not compared\\
$S^6$ and the Jacobian map $\to$ NS & Fluid coordinates, and gauge-scalar equations from the $S^6$ material carried to Navier--Stokes on $\mathbb T^3$; a Beltrami-mode forced solution from the $S^6$ oscillator & P & the forced solution verified exactly \cite{NSPaper}; the rest located (audit note \note{47\_}, edge 13)\\
$S^6\leftrightarrow$ ES & The number $5184=72^2$. It is $\det I$ in S6-5, where $I$ has Gram matrix $6G_{D_4}$, so $\det I=6^4\det D_4$; and it is $|\operatorname{disc}(A_5^4D_4)|=(\det A_5)^4\det D_4$ in S6-4 and R10 & A & both products $6^4\cdot4$ checked; no map between the two lattices was found\\
\bottomrule
\end{longtable}
\end{center}

Two further remarks. The Yang--Mills repository has an \texttt{s6/} topic entry, which holds the key-advances reader, and the record's frozen edition of 9 September was released together with companion Yang--Mills and Navier--Stokes editions (guide \fn{29\_}). The $S^6$ construction and the Jacobian counterexample are different objects, as the attribution file of the Yang--Mills repository states; the Jacobian counterexample is treated in its own paper \cite{JacobianPaper}. No crossing was found between the correction note and the zeta programme: the derived criterion (Theorem~\ref{thm:derived}) and the local algebra of Section~\ref{sec:correction} use nothing from the other workbenches.

\section{What remains open, and what was not checked}\label{sec:open}

\subsection{Open questions}

\begin{enumerate}[label=\arabic*.]
\item\label{q:construction} \emph{The construction and the counterexample.} Whether the constructed threefold exists with the stated properties, and so contradicts Theorem~2.2 of \cite{CDP20}, needs specialist verification. This is the record's open row \texttt{S6-OPEN-EXTERNAL-ADJUDICATION}; its other open row asks for independent replication of the lattice, divisor and residue computations (Section~\ref{sec:ledger}). Engel's exposition \cite{EngelS6} is a further independent source, and its comparison with the record has not been made.
\item\label{q:step} \emph{The published step.} The source manuscript (its §10.1) quotes the published sentence verbatim: from the conormal sequence, provided $L|_S$ is non-torsion, it suffices to prove the vanishing of $H^0(S,\widetilde\Omega^1_S\otimes L|_S)$. Read so, the step is the implication of Section~\ref{sec:correction} with $A=L|_S$ non-torsion. The note's Lemma~4.3 shows that the test family's $A$ can be chosen of this form on the local total space. The paper itself was not re-read here. By Corollary~\ref{cor:normal}, the implication can fail only on non-normal fibres.
\item\label{q:global} \emph{The global question left by the note.} The note leaves open whether the global hypotheses of \cite{CDP20} force $\delta_A$ to be injective on non-normal fibres. The record's answer at its own candidate is negative: at the cusp fibre the group $H^0(W,\Omega^1_X|_W\otimes A)$ is nonzero for every $L$ (the manuscript's Lemma~10.3 and Theorem~10.5; the monograph's Theorems~13.5 and~13.25). By Proposition~\ref{prop:kernel}, at a reduced fibre with normal crossings the question is exactly whether $H^0(\tilde S,\eta^*A)=0$. In the test family $\delta_A=0$ because $H^1(S,A)=0$ (Proposition~\ref{prop:sharp}).
\item\label{q:twoforms} \emph{Relative two-forms at the cusp.} The torsion of $\Omega^2_{X/B}$ on the record's charts is computed in Proposition~\ref{prop:torsion}(d): along the double curves it is the structure sheaf of the curve, and at a triple point the canonical module of the three branches. Its global form along the cusp fibre, and its effect on the direct images $R^qf_*\Omega^p_X$ through which the manuscript's Remark~1.3 would compute the Hodge numbers, were not examined. The record computes $h^{1,2}=1$ by another route (its ledger row~17).
\item\label{q:deformation} \emph{The deformation of S6-2.} The record does not claim that different values of the parameter give non-isomorphic complex manifolds, or that the family exhausts the deformations (key-advances reader, item~2).
\item\label{q:threefold} \emph{A threefold without the candidate.} Is there a compact complex threefold with the following properties?
\begin{itemize}
\item $H^1(X;\ZZ)=H^2(X;\ZZ)=0$, $a(X)=1$, a holomorphic algebraic reduction, and $b_3\ge2$;
\item a reduced fibre with normal crossings whose normalization has a component with $H^1(\mathcal O)=0$.
\end{itemize}
By Proposition~\ref{prop:kernel} and the argument of the manuscript's Theorem~10.5, such an $X$ would fail the relative hypothesis of \cite{CDP20} for every $L$, independently of the $S^6$ construction. This question was formulated in the first referee pass.
\item\label{q:hodge} \emph{Hodge numbers through relative forms.} Does the route of the manuscript's Remark~1.3, through $R^qf_*\Omega^p_{X/B}$ and the torsion computed in Proposition~\ref{prop:torsion}, reproduce $h^{1,2}=1$? (Also formulated in the first referee pass.)
\item\label{q:glue} \emph{The two glue codes.} The Erdős--Straus workbench's order-$72$ glue code for $A_5^4D_4$ and the one used in S6-4 have not been compared (Section~\ref{sec:overlaps}).
\item\label{q:paramodular} \emph{The paramodular character.} Is $e(\Phi/12)$ the restriction to $G_{\ZZ}$ of the character $\chi_3\chi_4^{-1}$ of Gritsenko and Hulek \cite[Lemma~1.1]{GritsenkoHulek}? The two characters have the same order and the same restrictions to the Levi group, the standard Heisenberg group and the centre of a Jacobi parabolic (Section~\ref{ssec:theories}). What remains is the value of $\chi_3\chi_4^{-1}$ on the minimal positive transvections $L_T$ and $t_\gamma$, which lie at isotropic vectors of divisors $6$ and $1$. If the two values are equal and of order $12$, then $G=G_{\ZZ}\cap\ker\chi_3\chi_4^{-1}$; if they are inverse to each other, the kernel is the twist $G_{-1}$ of Proposition~\ref{prop:twist}.
\item\label{q:arfone} \emph{The Arf-one torus.} Is the first homology of the Lie-framed torus of the higher-torus update (\fn{15\_}, (147)), whose spin form has Arf invariant one, the lattice $\langle\bar u,\bar w\rangle$ of Proposition~\ref{prop:heis}? If so, the odd spin structure that cuts out the Heisenberg part of the monodromy group is realised geometrically in the construction.
\item\label{q:elliptic} \emph{The modular part.} Is the $SL_2(\ZZ)$-graded piece $\gamma^\perp/\langle\gamma\rangle$ the homology of an elliptic family inside or under the tori?
\item\label{q:index12} \emph{The number $12$.} The index $12$ of $G$ in $G_{\ZZ}$ equals the order of $SL_2(\ZZ)^{\mathrm{ab}}$ and of the peripheral character. Theorem~\ref{thm:periph} explains the coincidence at the level of groups. Does it have a meaning in the geometry of the completed threefold, for instance through the mod-$24$ boundary of the update (\fn{15\_}, Theorem~6.1)?
\end{enumerate}

\subsection{What was not checked here}

\begin{itemize}
\item \emph{The construction.} Sections~6--7 of the monograph (the period-family descent, the fillings, the gluing, the nearby-cycle and Leray comparisons) were not read. The topology and sphere recognition, the algebraic dimension, and the Hodge, Bott--Chern and Aeppli computations were not read either. Every statement about them in this paper is the record's.
\item \emph{The monograph beyond the parts named.} Of its 154 pages, Part~I, the ledger, Theorem~13.6 and §§13.11 and~14 were read. The rest was mapped from its table of contents.
\item \emph{The period system beyond its integral monodromy.} Section~\ref{sec:monodromy} uses only the matrices $T_1,T_2$ and the form $Q_0$ of the manuscript's §2. The period functions, the Hodge filtration and the fillings were not examined.
\item \emph{The later papers and the archive.} The two papers of the topology supplement (\fn{15\_}, \fn{16\_}) were read only through their guide and their front matter. The proofs of S6-1 to S6-4 are in the archive \fn{28\_}, which was not downloaded.
\item \emph{The conversation and workflow files} \fn{07\_}--\fn{13\_} were not used.
\item \emph{The literature.} The paper \cite{CDP20} was not re-read; its statements are cited from the record and the note. No novelty search was made for the additions of Section~\ref{sec:correction}. The normal case of Corollary~\ref{cor:normal} is the source manuscript's Remark~10.4, and the general facts used are cited where they are used.
\end{itemize}

\appendix
\section{Catalogue of statements and files}\label{app:catalogue}

\subsection{Statements}

Every statement of the record that this paper states, proves, reproduces or locates, with where it is in the record and how far it was checked, and the new results. The labels are those defined in the Introduction. File numbers refer to the record; ``M'' is the monograph \fn{01\_} (``item'' is an item of its status list), ``S'' the reconstructed source manuscript \fn{03\_}, and ``N'' the correction note \fn{00\_}. Identifiers in typewriter type are rows of the claim ledger in \fn{04\_}.

\begin{center}\small
\begin{longtable}{@{}>{\raggedright\arraybackslash}p{0.36\linewidth}>{\raggedright\arraybackslash}p{0.21\linewidth}>{\raggedright\arraybackslash}p{0.24\linewidth}>{\raggedright\arraybackslash}p{0.11\linewidth}@{}}
\toprule
Statement & Where in the record & Status here & Here\\
\midrule
\endhead
\multicolumn{4}{@{}l}{\emph{The construction and the counterexample claim (record: \textsf{VERIFIED})}}\\*
$X\to\Pp^1$ compact complex, $X\cong_{\mathrm{diff}}S^6$, $a(X)=1$, $c_3(X)=2$ & S, summary; M §§6--10 & Located & §§\ref{sec:candidate}, \ref{sec:ledger}\\
$\pi_1(X)\cong\ZZ/|12\ell_0-4\ell_1-3\ell_2|$; the record's triple $(0,1,-1)$ gives $-1$ & S §7 & Located; the arithmetic reproduced & §\ref{sec:candidate}\\
Counterexample to Theorem~2.2 of \cite{CDP20} & M §13.11, item 5 & Located & §\ref{sec:ledger}\\
Hodge diamond, Frölicher sequence; Bott--Chern and Aeppli numbers; global representatives & M, status list items 17, 19, 20 & Located & §\ref{sec:ledger}\\
\midrule
\multicolumn{4}{@{}l}{\emph{The four rows labelled \textsf{REFUTED}}}\\*
The singular-fibre inference of \cite{CDP20}, p.~692 (\texttt{S6-CDP-P692}) & M, item 7; §13.11, item 1; N & Test family audited; see below & §\ref{sec:correction}\\
The direct-product formula of \cite{CDP98}, Lemma 3.2 (\texttt{S6-CDP98-MONO}) & M, item 12; §13.11, item 2 & Located & §\ref{sec:ledger}\\
Torsion of $\Omega^1_{X/B}$ at all three singular fibres (S, Remarks 1.3 and 9.21; \texttt{S6-HODGE-CUSP-TORSION}) & M, item 18; §14, item 6 & Proved here: correct for $p=1$; for $p=2$ the torsion is present & Prop.~\ref{prop:torsion}\\
Remark 7.28 for a fixed triple: $(100,1,-1)$ gives $1199$ (\texttt{S6-LOCAL-728}) & M §14, item 2 & Reproduced & §\ref{sec:ledger}\\
\midrule
\multicolumn{4}{@{}l}{\emph{The correction note}}\\*
$T_S$ is the torsion of $\Omega^1_S$ & N, Lemma 1.1 & Audited & §\ref{sec:correction}\\
Filtration and exact criterion & N, Lemma 2.1, Props.\ 2.2--2.3 & Audited; valid in any dimension and for vector bundles & Prop.~\ref{prop:filtration}\\
Normalization--conductor sequence and conductor differential & N, Lemma 3.1, Prop.\ 3.2, Cor.\ 3.3; S, Lemma 10.3 & Audited & §\ref{sec:correction}\\
The nodal-cubic test family & N, Prop.\ 4.1, Lemmas 4.2--4.3, Thm.\ 4.4; M, Thm.\ 13.6 & Audited; multiplier $1/3$ reproduced & Thm.~\ref{thm:testfamily}\\
Derived global criterion; relative form & N, Thm.\ 5.2, Cor.\ 5.3; M, Thm.\ 13.28 & Audited; generalized to $b_1=b_2=0$ with any nonzero effective divisor & Thm.~\ref{thm:derived}\\
Local relative differentials ($t=xy$, $xyz$, $s^m$) & N, Prop.\ 6.1 & Audited; general form proved here & Prop.~\ref{prop:torsion}\\
\midrule
\multicolumn{4}{@{}l}{\emph{New in version~2}}\\*
Torsion of a hypersurface's differentials: $T_{S,x}\cong\operatorname{Ext}^1(R/J,R)$, zero iff $S$ is normal at $x$; its global form & --- & Proved here; classical in general form (Vetter, Lebelt) & Prop.~\ref{prop:normal}\\
The step is valid on normal fibres & S, Remark 10.4 & Located in the record; the converse proved here & Cor.~\ref{cor:normal}\\
The kernel sheaf $K_S\cong\mathcal{H}om(\mathcal J,\mathcal O_S)\otimes N^*$; at a normal-crossing fibre the ambient group is $H^0(\tilde S,\eta^*(N^*\otimes A))$ & --- & Proved here (found in the first referee pass) & Prop.~\ref{prop:kernel}\\
Test family: $T_S\cong\mathcal O_D$, $H^1(S,A)=0$, $\delta_A=0$, ambient group $\cong\C$ & --- & Proved here & Prop.~\ref{prop:sharp}\\
Torsion of $\Omega^1_{R/\C\{t\}}$ is $R/(\gcd\partial_it)$, globally $(f^*\omega_B\otimes\mathcal O(Z))|_Z$; the torsion of relative two-forms at $t=xy$, $xyz$ (the canonical module of the three axes), $s^m$ & --- & Proved here (the global form found in the first referee pass) & Prop.~\ref{prop:torsion}\\
The arithmetic of S6-5, the formula for $\tau$ on $I$, and the exact value set & 26, item 5 & Proved here; (b)--(c) also in \cite[Prop.~4.7]{WorkbenchReader}; (d) found in the first referee pass & Prop.~\ref{prop:s65}\\
\midrule
\multicolumn{4}{@{}l}{\emph{New in version~3: the integral monodromy of the period system}}\\*
The period system: $T_1^3=T_2^4=I$, $T_0=I+N$, $N^2=0$; the invariant forms are $\ZZ Q_0$, of type $(1,6)$ & S, §2, Lemmas 2.2 and 2.8 & Reproduced & §\ref{sec:monodromy}\\
$G_{\ZZ}=SL_2(\ZZ)\ltimes\mathrm{Heis}(\ZZ)$ & --- & Proved here & Lemma~\ref{lem:stab}\\
$G=\ker\Phi$, index $12$, cyclic quotient; Levi image $SL_2(\ZZ)$; level exactly $12$; arithmetic and Zariski dense & --- & Proved here & Thm.~\ref{thm:mono}\\
The Heisenberg part: the odd refinement, the subgroups $N_k$, $v_H$ & --- & Proved here & Prop.~\ref{prop:heis}\\
The twelve twists; the extension does not split & --- & Proved here & Prop.~\ref{prop:twist}\\
The peripheral character: $\chi=-\operatorname{ab}$, the triple $(-1,4,-3)$, $\chi(M)\equiv s-6(p^2+pr+r^2)$ on $G$ & 15, §3.2, Thm.\ 6.1; 14, §2 & Proved here & Thm.~\ref{thm:periph}\\
The orbits of the $(3,4,\infty)$ covers at every level & --- & Proved here & Thm.~\ref{thm:orbitsS6}, Cor.~\ref{cor:orbitsS6}\\
No polarised limit mixed Hodge structure at the cusp: $b_N=\operatorname{diag}(6,-1)$ & S, Remark 3.23 & Reproduced & §\ref{ssec:mononeg}\\
$Sp(Q_0,\ZZ)$ is the paramodular group of level $6$; the restrictions of $e(\Phi/12)$ & --- & Proved here; the comparison with \cite{GritsenkoHulek} is Question~\ref{q:paramodular} & §\ref{ssec:theories}\\
\midrule
\multicolumn{4}{@{}l}{\emph{The topology supplement (5 September)}}\\*
The angle-forgetting map is null-homotopic & 14; 15 & Stated & §\ref{sec:later}\\
$[H\circ d]=\eta_4\eta_5\neq0$ via Arf invariant $1$ & 14; 15, Thm.\ 6.8, Cor.\ 6.12 & Arf invariant reproduced; the rest stated & §\ref{sec:later}\\
Hopf morphism, $12\nu_6$; peripheral morphism $\Psi_H$ & 15, Thms.\ 6.13, 6.15; 16, Thms.\ 11.1, 12.1, 12.3 & Values and kernel reproduced; classical input cited & §\ref{sec:later}\\
Local higher-torus map; Euler characteristic $6\neq2$ & 15 & Stated; arithmetic reproduced & §\ref{sec:later}\\
The octonionic cone $(B^{25},S^{24})$ & 15; 16 & Stated & §\ref{sec:later}\\
\midrule
\multicolumn{4}{@{}l}{\emph{The key advances (6 September)}}\\*
S6-1: the anticanonical ring $\C[U,V]$ & 26, item 1; archive & Stated; dimensions reproduced & §\ref{sec:later}\\
S6-2: a nonzero deformation & 26, item 2; archive & Stated; determinant ratio reproduced & §\ref{sec:later}\\
S6-3: covers of the finite fillings; attachment kernel $\ZZ$ & 26, item 3; archive & Stated & §\ref{sec:later}\\
S6-4: Niemeier--Leech arithmetic & 26, item 4; archive & Stated; $72$, $72$ roots, $81$ reproduced & §\ref{sec:later}\\
S6-5: the common lattice $I$ & 26, item 5; archive & Gram data reproduced; intersection stated & §\ref{sec:later}\\
\bottomrule
\end{longtable}
\end{center}

\subsection{Files}

\begin{center}\small
\begin{longtable}{@{}>{\raggedright\arraybackslash}p{0.2\linewidth}>{\raggedright\arraybackslash}p{0.5\linewidth}>{\raggedright\arraybackslash}p{0.22\linewidth}@{}}
\toprule
File & Content & Use here\\
\midrule
\endhead
\fn{00\_} (PDF, source zip) & The correction note & read in full; audited\\
\fn{01\_} & Verification monograph, 154 pp. & Part I, ledger, Thm.\ 13.6, §§13.11 and 14 read\\
\fn{02\_} & Annotated manuscript, 128 pp. & not read\\
\fn{03\_} & Clean reconstruction of the source manuscript, 108 pp. & summary, §2, Remarks 1.3, 3.23 and 9.21, §§10.1--10.3 read\\
\fn{04\_} & Sources and evidence, 20.8~MB zip & its claim ledger read\\
\fn{05\_}, \fn{06\_} & Mathematical provenance and its receipt & \fn{05\_} read\\
\fn{07\_}--\fn{13\_} & Sanitised conversation and typed workflow annex, receipts, validation & not used\\
\fn{14\_} & Reader's guide to the topology supplement & read in full\\
\fn{15\_}, \fn{16\_} & The higher torus wrap (50 pp.); Real $S^{12}$ and $S^{24}$ routes (64 pp.) & front matter read\\
\fn{17\_}--\fn{19\_} & Their sources, notes and manifest & \fn{18\_} read in part\\
\fn{26\_}, \fn{27\_} & Key-advances reader and its source & \fn{26\_} read in full\\
\fn{28\_} & Frozen project, 174~MB, 1{,}079 files & not downloaded\\
\fn{29\_}, \fn{30\_} & Guide and manifest of the frozen project & \fn{29\_} read\\
\texttt{README.md}; metadata, rights, licence and checksum files & Version 1.0 front file; record metadata & \texttt{README.md} read\\
\bottomrule
\end{longtable}
\end{center}

\section{How this paper was checked}\label{app:verification}

\paragraph{Sources.}
The record 10.5281/zenodo.22678442. Its file list and MD5 checksums, and the files added by each of its four versions, were read from the record's public metadata on 26 September 2026. Thirty-two of its 33 files were downloaded, with the owner's permission, and their MD5 checksums and sizes equal the record's; the 174~MB archive \fn{28\_} was not downloaded. The script \note{s6r\_record\_inventory.py} reproduces this. The parts of the record read, and those only mapped, are listed in Appendix~\ref{app:catalogue}.

\paragraph{What was refereed earlier.}
The overlap edges of Section~\ref{sec:overlaps} other than the one marked ``new'' come from the audit notes \note{28\_} and \note{47\_}. These were refereed (the ninth referee pass, and the referee pass on the notes \note{43\_}--\note{47\_}) and entered the paper on four workbenches \cite{WorkbenchReader}.

\paragraph{The first referee pass.}
On 27 September 2026 an independent Claude instance that had written none of the text refereed version~1 of 26 September. It read in both directions: for overstatement, and for understatement, missed generality and missed or implied results. It reran both scripts of version~1 and wrote its own checks, 22 of them, all passing; they are included as \note{referee/s6r\_referee\_checks.py}. It found no mathematical error. It made 2 major and 12 minor findings.
\begin{itemize}
\item \emph{Major findings.}
\begin{itemize}
\item The abstract listed as new that the step is valid on normal fibres, which is the source manuscript's Remark~10.4.
\item The record's two \textsf{OPEN} rows were misreported.
\end{itemize}
\item \emph{Minor findings.} Among them:
\begin{itemize}
\item the classical results of Vetter and Lebelt were not cited;
\item the phrase ``on any fibre'' should read ``on any reduced fibre'';
\item the example gluing triple was the paper's own, not the record's;
\item the description of the monograph's §14 was inaccurate;
\item Engel's citation is of the 1998 paper;
\item S6-1 to S6-4 were labelled ``located'' although their proofs were not seen.
\end{itemize}
\end{itemize}
Every finding was checked by the author against the sources before it was applied. For example, the claim ledger's two open rows were read in the evidence archive, and Engel's text and the bibliographic data of Vetter, Lebelt and Miller--Vassiliadou were read online. All findings were applied.

The pass also proposed missed or implied results, each with a proof. The author verified each proof, and each result is in this version with its credit:
\begin{itemize}
\item Proposition~\ref{prop:kernel};
\item the global form of the torsion of relative one-forms;
\item the extension of Theorem~\ref{thm:derived} to $b_1=b_2=0$, with its consequence for rational homology six-spheres;
\item the exact value set in Proposition~\ref{prop:s65}(d), checked numerically for $|n|\le3000$;
\item two open questions (Section~\ref{sec:open}).
\end{itemize}
Two further proposals were not adopted. A reflexivity criterion was left out as outside the paper's topic. A sketched global form of the torsion of relative two-forms was left out because only a sketch was given.

\paragraph{The referee pass of the monodromy computations.}
The results of Section~\ref{sec:monodromy} were first written in the audit note \note{52\_}. On 27 September 2026 an independent Claude instance that had written none of the note refereed it in both directions. It re-derived every proof, found no error that invalidates a stated theorem, and found that the first version said much less than its own data prove. The sharpened statements, among them Theorem~\ref{thm:mono} in its present form, Propositions~\ref{prop:heis} and~\ref{prop:twist}, and the orbit classification, come from that report. Each was re-derived by claude-ab and checked with code written separately from the referee's (\note{ver52\_referee\_claims.py}) before it was adopted. The referee's reading that $\Phi$ is the restriction of the character $\chi_3\chi_4^{-1}$ of Gritsenko and Hulek depends on a normalisation that was not read in full, and it is stated here as Question~\ref{q:paramodular}.

\paragraph{What rests on this author and one referee pass.}
\begin{itemize}
\item the map of the record (Sections~\ref{sec:record}, \ref{sec:ledger} and~\ref{sec:later});
\item the audit of the correction note (Section~\ref{sec:correction});
\item Lemma~\ref{lem:onerel}, Propositions~\ref{prop:normal}, \ref{prop:kernel}, \ref{prop:sharp}, \ref{prop:torsion} and~\ref{prop:s65}, and Corollary~\ref{cor:normal};
\item the generalizations of Theorem~\ref{thm:derived} and the remark on its citation;
\item the new overlap edge;
\item Section~\ref{sec:monodromy} and Section~\ref{ssec:theories}.
\end{itemize}

\paragraph{The scripts.}
They import no code from the record. They are in the folder \texttt{contrib/\allowbreak claude-ab/\allowbreak rewrite/\allowbreak checks/\allowbreak s6/} on the branch named below.
\begin{center}\small
\begin{tabularx}{\linewidth}{@{}>{\raggedright\arraybackslash}p{0.27\linewidth}>{\raggedright\arraybackslash}X>{\raggedright\arraybackslash}p{0.13\linewidth}@{}}
\toprule
Script & What it checks & Time, result\\
\midrule
\note{s6r\_checks.py} &
\begin{minipage}[t]{\linewidth}\raggedright
47 symbolic and exact checks, in four groups:
\begin{itemize}[leftmargin=1em]
\item \emph{the test family}: smoothness of the total space, the discriminant, the normalization and the local model at the node, the sections $O,P$ and the multiplier $1/3$, $H^0(A_\lambda)=0$, the relations $u\tau=v\tau=0$, the map $\mu$ of Proposition~\ref{prop:sharp}, and the independence of $dU,dV$;
\item \emph{the local modules}: the gcd formula on nine monomials with unit factors, the torsion of relative one- and two-forms at $t=xy$, $xyz$, $s^m$, and the presentation of the torsion at $t=xyz$;
\item \emph{the derived criterion}: the Riemann--Roch coefficient $\tfrac12$ of $c_3$;
\item \emph{the finite statements}: $1199$; the Arf invariant; $\Psi_H$; the Euler characteristic; $\det_{\R}\Pi=\operatorname{Im}\tau\cdot D$; the numbers of S6-1, S6-2, S6-4 and S6-5, including the quaternion formula for $\tau$.
\end{itemize}
\end{minipage} & 2 s; all pass\\
\note{s6r\_record\_inventory.py} & the 33 files of the record against the downloaded copies, by MD5 and size & under 1 s; 32 match, one not local\\
\note{referee/s6r\_referee\_checks.py} & the first referee's own checks: the kernel sheaf at a triple point, Hilbert functions of the relative forms at $t=xy$, $xyz$, $x^2y$ and two quadrics, the value set of S6-5, the source's matrices, the gluing triples and the claim-ledger counts & 7 s; 22 checks, all pass\\
\note{s6\_monodromy\_checks.py} & the relations $T_1^3=T_2^4=I$, $N^2=0$; the invariance of $Q_0$, its uniqueness up to scale and its elementary divisors; Lemma~\ref{lem:stab} symbolically; the generation of $SL_2(\ZZ)$ by $R$ and $S$; $c=h(1,0,-6)$; the refinements and the coset representatives; the indices of $G\bmod D$ for $D=2,3,4,6,8,9,12$ by breadth-first enumeration; the level-$12$ image; the orbits for $D\le9$ & 49 s; 24 checks, all pass\\
\note{ver52\_referee\_claims.py} & that $S\mapsto9$, $T\mapsto1$ defines a homomorphism on all $1152$ elements of $SL_2(\ZZ/12)$; the decompositions of Theorem~\ref{thm:mono}(b); $\Phi$ multiplicative on random products and zero on random words in $T_1,T_2$; $\Phi=0$ on all $165{,}888$ elements of $G\bmod12$; $H'=\langle c,c'\rangle$; the subgroups $N_k$; the non-splitting computation; the dual action; Theorem~\ref{thm:orbitsS6} for all $D\le14$ as full partitions; Theorem~\ref{thm:periph}; Proposition~\ref{prop:heis}(d) & 5 s; 22 checks, all pass\\
\note{referee/ref52\_partB\_group.py}, \note{referee/ref52\_connections.py} & the second referee's checks: a constructive membership algorithm for $G=\ker\Phi$, the orbits for $D\le24$, and the checks behind Section~\ref{ssec:theories}; its reruns of the two scripts above also pass & 14 s and under 1 s; 26 and 10 checks, all pass\\
\bottomrule
\end{tabularx}
\end{center}

\paragraph{What has not been done.}
The new results of versions~2 and~3, Proposition~\ref{prop:kernel} and Theorem~\ref{thm:mono} among them, have been checked by the author and by the one referee pass that proposed or read them, but not by a second independent pass. The paper \cite{CDP20} itself, the record's construction, and the archive \fn{28\_} were not read (Section~\ref{sec:open}).

\paragraph{Where everything is.}
This paper, its sources and its scripts are on the branch \texttt{claude/claude-ab-grind-20260925} of \texttt{KokunoYumeto/zeta-function-research-reader}: the paper in \texttt{contrib/\allowbreak claude-ab/\allowbreak rewrite/\allowbreak papers/\allowbreak s6/}, the scripts in \texttt{contrib/\allowbreak claude-ab/\allowbreak rewrite/\allowbreak checks/\allowbreak s6/} (the referees' in its folder \texttt{referee/}). Version~2 and the audit notes remain in the branch's history at commit \texttt{a90d9e5b}.

\section{Changes in version 3}\label{app:changes}

Version~2 of 27 September 2026 was withdrawn from the branch on the same day because of its framing, together with the audit notes. Version~3 keeps its theorems, proofs, citations and status lines, and makes the following changes.

\paragraph{Framing.}
\begin{itemize}
\item The front matter is a prose abstract and an introduction with the main results stated as Theorems~I--IV. The section \emph{About this reader} is replaced by the subsection \emph{How this record was made}, which gives the credit for the construction, the record and this paper, and the four statuses.
\item Remarks on how documents were produced, other than credit, are removed: the naming of the source on the monograph's title page, the correction note's author line, the record's self-description as a verification aid, the description of the record's labels as assigned by its verification, and the description of the conversation files. The statement that the source manuscript is unrefereed is removed.
\item The sentence that the paper takes no position on the construction is replaced by the statement that the construction and the counterexample claim are not verified here.
\item Section~\ref{sec:overlaps} is now \emph{Bridges to other areas and to the author's other workbenches}, and treats the connections to established theories first. The remark on the name of the construction is reduced to the statement that the $S^6$ construction and the Jacobian counterexample are different objects.
\item The paper is called a paper, and the companion texts are cited in their current versions.
\end{itemize}

\paragraph{New results.}
\begin{itemize}
\item Section~\ref{sec:monodromy}, the integral monodromy of the period system: Lemma~\ref{lem:stab}, Theorem~\ref{thm:mono}, Propositions~\ref{prop:heis} and~\ref{prop:twist}, Theorem~\ref{thm:periph}, Theorem~\ref{thm:orbitsS6} and Corollary~\ref{cor:orbitsS6}, and the proved negative statements of Section~\ref{ssec:mononeg}. They were first written in the audit note \note{52\_} and refereed there (Appendix~\ref{app:verification}).
\item Section~\ref{ssec:theories}: the paramodular group of level $6$ and the restrictions of $e(\Phi/12)$, spin structures, the Jordan type at the cusp, the modular part, and arithmeticity.
\item Questions~\ref{q:paramodular}--\ref{q:index12}.
\end{itemize}

\paragraph{Corrections.}
\begin{itemize}
\item The edge $S^6\to$ ES now states that the orbits of the covers at every level are determined (Theorem~\ref{thm:orbitsS6}); version~2 stated only that the package's statements are correct for odd levels and do not hold as stated at even levels, which remains true.
\end{itemize}

\addcontentsline{toc}{section}{References}
\begin{thebibliography}{99}
\small\raggedright
\bibitem{AlpogeS6} \emph{A compact complex threefold fibred by tori over the projective line, and the six-sphere}, manuscript circulated by L.~Alpöge, 2026, \url{https://alpo.ge/s6.pdf}. Its opening summary is titled \emph{The $(3,4,\infty)$ modular family of 2-tori, completed at its three special points, is a complex structure on $S^6$}.
\bibitem{AS68} M.~F.~Atiyah, I.~M.~Singer, \emph{The index of elliptic operators: III}, Ann. of Math. (2) \textbf{87} (1968), 546--604.
\bibitem{BH93} W.~Bruns, J.~Herzog, \emph{Cohen--Macaulay Rings}, Cambridge Studies in Advanced Mathematics \textbf{39}, Cambridge University Press, Cambridge, 1993.
\bibitem{CDP98} F.~Campana, J.-P.~Demailly, T.~Peternell, \emph{The algebraic dimension of compact complex threefolds with vanishing second Betti number}, Compositio Math. \textbf{112} (1998), 77--91.
\bibitem{CDP20} F.~Campana, J.-P.~Demailly, T.~Peternell, \emph{The algebraic dimension of compact complex threefolds with vanishing second Betti number}, Compositio Math. \textbf{156} (2020), 679--696; arXiv:1904.11179.
\bibitem{EngelS6} P.~Engel, \emph{Complex structures on $S^6$}, 13 September 2026, \url{https://philip-engel.github.io/S6.pdf}.
\bibitem{GR84} H.~Grauert, R.~Remmert, \emph{Coherent Analytic Sheaves}, Grundlehren der mathematischen Wissenschaften \textbf{265}, Springer, Berlin, 1984.
\bibitem{Har77} R.~Hartshorne, \emph{Algebraic Geometry}, Graduate Texts in Mathematics \textbf{52}, Springer, New York, 1977.
\bibitem{Hir66} F.~Hirzebruch, \emph{Topological Methods in Algebraic Geometry}, 3rd ed., Grundlehren der mathematischen Wissenschaften \textbf{131}, Springer, Berlin, 1966.
\bibitem{Mat86} H.~Matsumura, \emph{Commutative Ring Theory}, Cambridge Studies in Advanced Mathematics \textbf{8}, Cambridge University Press, Cambridge, 1986.
\bibitem{Toda62} H.~Toda, \emph{Composition Methods in Homotopy Groups of Spheres}, Annals of Mathematics Studies \textbf{49}, Princeton University Press, Princeton, 1962.
\bibitem{Wilson09} R.~A.~Wilson, \emph{Octonions and the Leech lattice}, J. Algebra \textbf{322} (2009), 2186--2190.
\bibitem{Leb74} K.~Lebelt, \emph{Torsion äußerer Potenzen von Moduln der homologischen Dimension 1}, Math. Ann. \textbf{211} (1974), 183--197.
\bibitem{MS74} J.~W.~Milnor, J.~D.~Stasheff, \emph{Characteristic Classes}, Annals of Mathematics Studies \textbf{76}, Princeton University Press, Princeton, 1974.
\bibitem{MV19} C.~Miller, S.~Vassiliadou, \emph{(Co)torsion of exterior powers of differentials over complete intersections}, J. Singul. \textbf{19} (2019), 131--162.
\bibitem{Ser56} J.-P.~Serre, \emph{Géométrie algébrique et géométrie analytique}, Ann. Inst. Fourier \textbf{6} (1956), 1--42.
\bibitem{Ser73} J.-P.~Serre, \emph{A Course in Arithmetic}, Graduate Texts in Mathematics \textbf{7}, Springer, New York, 1973.
\bibitem{Vet70} U.~Vetter, \emph{Äußere Potenzen von Differentialmoduln reduzierter vollständiger Durchschnitte}, Manuscripta Math. \textbf{2} (1970), 67--76.
\bibitem{SarnakThin} P.~Sarnak, \emph{Notes on thin matrix groups}, in: \emph{Thin Groups and Superstrong Approximation}, MSRI Publications \textbf{61}, Cambridge University Press, Cambridge, 2014, 343--362.
\bibitem{BravThomas} C.~Brav, H.~Thomas, \emph{Thin monodromy in $Sp(4)$}, Compositio Math. \textbf{150} (2014), 333--343; arXiv:1210.0523.
\bibitem{SerreTrees} J.-P.~Serre, \emph{Trees}, Springer, Berlin, 1980.
\bibitem{Johnson} D.~Johnson, \emph{Spin structures and quadratic forms on surfaces}, J. London Math. Soc. (2) \textbf{22} (1980), 365--373.
\bibitem{AtiyahSpin} M.~F.~Atiyah, \emph{Riemann surfaces and spin structures}, Ann. Sci. École Norm. Sup. (4) \textbf{4} (1971), 47--62.
\bibitem{GritsenkoHulek} V.~Gritsenko, K.~Hulek, \emph{Commutator coverings of Siegel threefolds}, Duke Math. J. \textbf{94} (1998), no.~3, DOI 10.1215/S0012-7094-98-09421-2; arXiv:alg-geom/9702007, whose numbering is used here.
\bibitem{Kulikov} V.~S.~Kulikov, \emph{Degenerations of K3 surfaces and Enriques surfaces}, Math. USSR Izv. \textbf{11} (1977), 957--989.
\bibitem{PerssonPinkham} U.~Persson, H.~Pinkham, \emph{Degeneration of surfaces with trivial canonical bundle}, Ann. of Math. (2) \textbf{113} (1981), 45--66.
\bibitem{ESPaper} Claude (Anthropic), \emph{The Erdős--Straus equation through residual shells}, version~4, 2026; source in \texttt{contrib/claude-ab/rewrite/papers/erdos\_straus/} on the branch \texttt{claude/\allowbreak claude-ab-grind-20260925} of \texttt{KokunoYumeto/\allowbreak zeta-function-research-reader}.
\bibitem{YMPaper} Claude (Anthropic), \emph{The Yang--Mills workbench: its results, with proofs, and its bridges}, version~3, 2026, in the Zenodo record 10.5281/zenodo.22883643 and its later versions; source in \texttt{contrib/claude-ab/rewrite/papers/yang\_mills/} on the same branch.
\bibitem{NSPaper} Claude (Anthropic), \emph{The Navier--Stokes workbench: the exact radius of the Rindler shear series, and curvature rates for the map to Yang--Mills}, 2026, in the Zenodo record 10.5281/zenodo.22883643 and its later versions; source in \texttt{contrib/claude-ab/rewrite/papers/navier\_stokes/} on the same branch.
\bibitem{JacobianPaper} Claude (Anthropic), \emph{The Jacobian-conjecture counterexample of Alpöge: fibres, a torus action, a Whitney-cusp normal form and the Weyl algebra}, 2026, in the same Zenodo record as this paper; its condensed version is in \texttt{contrib/claude-ab/rewrite/condensed/} on the same branch.
\bibitem{WorkbenchReader} Claude (Anthropic), \emph{Four workbenches and where they meet: a reader of the Erdős--Straus, Collatz, Erdős Problem 817 and Yang--Mills workbenches}, 2026; withdrawn from the branch on 27 September 2026 and kept in its history at commit \texttt{a90d9e5b}, folder \texttt{contrib/claude-ab/grind\_20260925/workbench\_reader/}.
\end{thebibliography}

\end{document}
